\documentclass[a4paper, 12pt, oneside, openany]{article}
\frenchspacing
\usepackage[left=3cm,right=1.5cm,top=2cm,bottom=2cm]{geometry}
\pagestyle{plain}

% ✅ Кодировка и шрифты (ПРАВИЛЬНЫЙ ПОРЯДОК)
\usepackage[T2A]{fontenc}      % ✅ Только T2A для русского!
\usepackage[utf8]{inputenc}
\usepackage{lmodern}           % ✅ Масштабируемые шрифты
\usepackage[russian]{babel}

% ✅ Остальные пакеты
\usepackage{float}
\usepackage{flafter}
\usepackage{wrapfig}
\usepackage[pdftex,unicode]{hyperref}
% ... гиперссылки настройки ...

% ✅ Microtype (после шрифтов, с final)
\usepackage[expansion=false, protrusion=false]{microtype}  % ✅ Отключить всё проблемное

\usepackage{amsmath, amssymb}
\usepackage[warn]{mathtext}

\usepackage{pgfplots}
\pgfplotsset{compat=newest}
\usepackage{tikz}              % ✅ ДОБАВИТЬ СЮДА
\usetikzlibrary{shapes,arrows,positioning,calc}


%% Реализация библиографии пакетами biblatex и biblatex-gost с использованием движка biber - не работает в overleaf из-за долгой компиляции в бесплатном режиме. В других приложениях можно брать %%%
%
%\usepackage{csquotes} % Пакет для оформления сложных блоков цитирования (biblatex рекомендует его подключать)
%\usepackage[%
%backend=biber,                                      % Движок
%bibencoding=utf8,                                   % Кодировка bib-файла
%sorting=none,                                       % Настройка сортировки списка литературы
%style=gost-numeric,                                 % Стиль цитирования и библиографии по ГОСТ
%language=auto,                                      % Язык для каждой библиографической записи задается отдельно
%autolang=other,                                     % Поддержка многоязычной библиографии
%sortcites=true,                                     % Если в квадратных скобках несколько ссылок, то отображаться будут отсортированно
%movenames=false,                                    % Не перемещать имена, они всегда в начале библиографической записи
% maxnames=5,                                         % Максимальное отображаемое число авторов
% minnames=3,                                         % До скольки сокращать число авторов, если их больше максимума
% doi=false,                                          % Не отображать ссылки на DOI
% isbn=false,                                         % Не показывать ISBN, ISSN, ISRN
% ]{biblatex}[2016/09/17]
% \DeclareDelimFormat{bibinitdelim}{}                 % Убираем пробел между инициалами (Иванов И.И. вместо Иванов И. И.)
% \addbibresource{biba.bib}                           % Определяем файл с библиографией
%
% %%% Скрипт, который автоматически подбирает язык (и, следовательно, формат) для каждой библиографической записи %%%
% %%% Если в названии работы есть кириллица - меняем значение поля langid на russian %%%
%%%% Все оставшиеся пустые места в поле langid заменяем на english %%%
%
%\DeclareSourcemap{
%   \maps[datatype=bibtex]{
%     \map{
%         \step[fieldsource=title, match=\regexp{^\P{Cyrillic}*\p{Cyrillic}.*}, final]
%         \step[fieldset=langid, fieldvalue={russian}]
%     }
%     \map{
%         \step[fieldset=langid, fieldvalue={english}]
%     }
%   }
% }


%%% Задаем параметры оформления рисунков и таблиц %%%

\usepackage{graphicx, caption, subcaption}  % Подгружаем пакеты для работы с графикой и настройки подписей
\graphicspath{{images/}} % Определяем папку с рисунками
\RequirePackage{caption}
\DeclareCaptionLabelSeparator{defffis}{ - }
\captionsetup[figure]{font=small,  name=Рисунок, justification=centering, labelsep=defffis, labelfont=bf} % Задаем параметры подписей к рисункам: маленький шрифт (в данном случае 12pt), ширина равна ширине текста, полнотекстовая надпись "Рисунок", выравнивание по центру
\captionsetup[subfigure]{font=small} % Индексы подрисунков а), б) и так далее тоже шрифтом 12pt (по умолчанию делает еще меньше)
\captionsetup[table]{singlelinecheck=false,font=small,width=\textwidth,justification=justified, labelsep=defffis, labelfont=bf} % Задаем параметры подписей к таблицам: запрещаем переносы, маленький шрифт (в данном случае 12pt), ширина равна ширине текста, выравнивание по ширине
% \captiondelim{-} % Разделителем между номером рисунка/таблицы и текстом в подписи является длинное тире. Не работает
\setkeys{Gin}{width=\textwidth} % По умолчанию размер всех добавляемых рисунков будет подгоняться под ширину текста
\renewcommand{\thesubfigure}{\asbuk{subfigure}} % Нумерация подрисунков строчными буквами кириллицы
%\setlength{\abovecaptionskip}{0pt} % Отбивка над подписью - тут не меняем
%\setlength{\belowcaptionskip}{0pt} % Отбивка под подписью - тут не меняем
\usepackage[section]{placeins} % Объекты типа float (рисунки/таблицы) не вылезают за границы секциии, в которой они объявлены

\usepackage{longtable} %можно делать многостраничные таблицы
\usepackage{rotating} %можно поворачивать объекты на заданный угол \rotatebox{90}{}
\usepackage{multirow}
\usepackage{color,colortbl}
\definecolor{darkishgreen}{RGB}{39,203,22}
\definecolor{LightCyan}{rgb}{0.88,1,1}
\definecolor{Gray}{gray}{0.9}
\definecolor{lightRed}{RGB}{230,170,150}
\definecolor{modRed}{RGB}{230,82,90}
\definecolor{strongRed}{RGB}{230,6,6}
\setlength{\parindent}{1.25cm} % Устанавливает отступ первой строки 1.25 см
\usepackage{enumitem}
\renewcommand*{\labelitemi}{\normalfont{-}}        % В ненумерованных списках для пунктов используем короткое тире
\makeatletter
    \AddEnumerateCounter{\asbuk}{\russian@alph}     % Объясняем пакету enumitem, как использовать asbuk
\makeatother
\renewcommand{\labelenumii}{\asbuk{enumii})}        % Кириллица для второго уровня нумерации
\renewcommand{\labelenumiii}{\arabic{enumiii})}     % Арабские цифры для третьего уровня нумерации
\setlist{noitemsep, leftmargin=*}                   % Убираем интервалы между пунками одного уровня в списке
% \setlist[1]{labelindent=\parindent}                 % Отступ у пунктов списка равен абзацному отступу
% \setlist[2]{leftmargin=\parindent}                  % Плюс еще один такой же отступ для следующего уровня
% \setlist[3]{leftmargin=\parindent}                  % И еще один для третьего уровня

\setlist[itemize]{leftmargin=-0.15cm, itemindent=1.25cm} % Для всех itemize
\setlist[enumerate]{leftmargin=0.25cm, itemindent=1.25cm} % Для всех enumerate
\usepackage{pdflscape}
\linespread{1.25}
\clubpenalty=10000                                  % Запрещаем разрыв страницы после первой строки абзаца
\widowpenalty=10000                                 % Запрещаем разрыв страницы после последней строки абзаца
\brokenpenalty=4991                                 % Ограничение на разрыв страницы, если строка заканчивается переносом

\usepackage{appendix} % Для удобного создания приложений
\usepackage{tocloft}
\renewcommand{\cftsecleader}{\cftdotfill{\cftdotsep}}
% \newcommand{\intro}[1]{
%     \stepcounter{section}
%     \section*{\centeringПриложение \Asbuk{section}}
%     \begin{center}
%         \bf{#1}
%     \end{center}
%     \markboth{\MakeUppercase{#1}}{}
%     \addcontentsline{toc}{section}{Приложение \Asbuk{section}. #1}
% }
\usepackage{mdframed}
\usepackage{xcolor} % Для цвета рамки (опционально)
\usepackage{listings}
\DeclareCaptionFont{white}{\color{white}} %% это сделает текст заголовка белым
%% код ниже нарисует серую рамочку вокруг заголовка кода.
\DeclareCaptionFormat{listing}{\colorbox{gray}{\parbox{\textwidth}{#1#2#3}}}
\captionsetup[lstlisting]{format=listing,labelfont=white,textfont=white, singlelinecheck=false,font=small,width=\textwidth,justification=justified, labelsep=defffis}

\lstset{ 
language=Python,
basicstyle=\ttfamily\small, % размер и начертание шрифта для подсветки кода
numbers=left,               % где поставить нумерацию строк (слева\справа)
numberstyle=\footnotesize,           % размер шрифта для номеров строк
stepnumber=1,                   % размер шага между двумя номерами строк
numbersep=5pt,                % как далеко отстоят номера строк от подсвечиваемого кода
showspaces=false,            % показывать или нет пробелы специальными отступами
showstringspaces=false,      % показывать или нет пробелы в строках
showtabs=false,             % показывать или нет табуляцию в строках
tabsize=2,                 % размер табуляции по умолчанию равен 2 пробелам
captionpos=t,              % позиция заголовка вверху [t] или внизу [b] 
breaklines=true,           % автоматически переносить строки (да\нет)
breakatwhitespace=false, % переносить строки только если есть пробел
% escapeinside={\%}{*)},   % если нужно добавить комментарии в коде
escapeinside={(*@}{@*)},  % ✅ Разные открывающий и закрывающий
columns=fullflexible,
flexiblecolumns,
extendedchars=true,
inputencoding=utf8,
  % Для кириллицы в комментариях:
  literate={а}{{\cyra}}1 {б}{{\cyrb}}1 {в}{{\cyrv}}1 {г}{{\cyrg}}1 
           {д}{{\cyrd}}1 {е}{{\cyre}}1 {ё}{{\cyryo}}1 {ж}{{\cyrzh}}1
           {з}{{\cyrz}}1 {и}{{\cyri}}1 {й}{{\cyrishrt}}1 {к}{{\cyrk}}1
           {л}{{\cyrl}}1 {м}{{\cyrm}}1 {н}{{\cyrn}}1 {о}{{\cyro}}1
           {п}{{\cyrp}}1 {р}{{\cyrr}}1 {с}{{\cyrs}}1 {т}{{\cyrt}}1
           {у}{{\cyru}}1 {ф}{{\cyrf}}1 {х}{{\cyrh}}1 {ц}{{\cyrc}}1
           {ч}{{\cyrch}}1 {ш}{{\cyrsh}}1 {щ}{{\cyrshch}}1 {ъ}{{\cyrhrds}}1
           {ы}{{\cyrery}}1 {ь}{{\cyrsftsn}}1 {э}{{\cyre}}1 {ю}{{\cyryu}}1
           {я}{{\cyrya}}1
           {А}{{\CyrA}}1 {Б}{{\CyrB}}1 {В}{{\CyrV}}1 {Г}{{\CyrG}}1
           {Д}{{\CyrD}}1 {Е}{{\CyrE}}1 {Ё}{{\CyrYo}}1 {Ж}{{\CyrZh}}1
           {З}{{\CyrZ}}1 {И}{{\CyrI}}1 {Й}{{\CyrIi}}1 {К}{{\CyrK}}1
           {Л}{{\CyrL}}1 {М}{{\CyrM}}1 {Н}{{\CyrN}}1 {О}{{\CyrO}}1
           {П}{{\CyrP}}1 {Р}{{\CyrR}}1 {С}{{\CyrS}}1 {Т}{{\CyrT}}1
           {У}{{\CyrU}}1 {Ф}{{\CyrF}}1 {Х}{{\CyrH}}1 {Ц}{{\CyrTs}}1
           {Ч}{{\CyrCh}}1 {Ш}{{\CyrSh}}1 {Щ}{{\CyrShch}}1 {Ъ}{{\CyrHrds}}1
           {Ы}{{\CyrEry}}1 {Ь}{{\CyrSftsn}}1 {Э}{{\CyrE}}1 {Ю}{{\CyrYu}}1
           {Я}{{\CyrYa}}1,
keywordstyle=\color{blue}\bfseries,
identifierstyle=\color{black},
morecomment=[l]{\#},  % ✅ Экранировать символ #
commentstyle=\color{green}\ttfamily,
stringstyle=\color{red}\ttfamily,
morestring=[b]",
keepspaces=true,
% escapechar=\%,
texcl=false,
backgroundcolor=\color{white}, % цвет фона подсветки - используем \usepackage{color}
frame=single              % рисовать рамку вокруг кода
}
% \usepackage{minted}
% \usemintedstyle{colorful}
% \DeclareCaptionFont{white}{\color{white}} %% это сделает текст заголовка белым
% %% код ниже нарисует серую рамочку вокруг заголовка кода.
% \DeclareCaptionFormat{minted}{\colorbox{gray}{\parbox{\textwidth}{#1#2#3}}}
% \captionsetup[minted]{format=minted,labelfont=white,textfont=white, singlelinecheck=false,font=small,width=\textwidth,justification=justified, labelsep=defffis}


\begin{document}
\begin{titlepage} % начало титульной страницы
\begin{center} % включить выравнивание по центру
Федеральное государственно автономное образовательное учреждение\\ высшего образования\\
<<Московский физико-технический институт (государственный университет)>>\\
Центр дополнительного, дополнительного профессионального\\и онлайн-образования <<Пуск>>\\ [2cm]
\end{center} % закончить выравнивание по центру
\par
\textbf{Направление подготовки:} 01.04.02 Прикладная математика и информатика

\textbf{Направленность (профиль) подготовки:} Науки о данных

\vspace{2cm}

\begin{center} % включить выравнивание по центру
{\huge\textbf{Разработка функции потерь для обеспечения эмоциональной согласованности между речью и мимикой при дообучении аудио-управляемых моделей генерации говорящих лиц}}\\
(магистерская диссертация)\\[2cm]
\end{center} % закончить выравнивание по центру
\flushright {
\begin{minipage}{0.5\textwidth} % начало маленькой врезки в половину ширины текста
\begin{flushleft} % выровнять её содержимое по левому краю
\textbf{Студент:}\\
Почтар Катрин Викторовна\\
\hrulefill
\parskip=-5mm
\begin{center}\footnotesize\textit{(подпись студента)}\end{center}
\parskip=3mm

\textbf{Научный руководитель:}
\par
Садеков Тимур Шамилевич\\
\hrulefill
\parskip=-5mm
\begin{center} \footnotesize\emph{(подпись научного руководителя)}\end{center}
\parskip=3mm

\end{flushleft} % конец выравнивания по левому краю
\end{minipage} % конец врезки
}
\vfill % заполнить всё доступное ниже пространство
\begin{center}{Москва \the\year}\end{center} % вывести дату
\thispagestyle{empty} % не нумеровать страницу
\end{titlepage} % конец титульной страницы

\newpage
\setcounter{page}{2}
\section{Аннотация}

Диссертация посвящена разработке кросс-модальной функции потерь, повышающей эмоциональную согласованность речи и мимики при дообучении архитектур генерации говорящих лиц Wav2Lip и SadTalker на четырёх эмоциях RAVDESS (happy, sad, angry, disgust). Существующие эмоциональные подходы требуют замены архитектуры или меток на инференсе.

Предложена параметризуемая функция потерь, надстраиваемая над L1-реконструкцией и связывающая логиты двух зафиксированных энкодеров --- wav2vec~2.0 для аудио и TimeSformer для видео. По итогам пятиконфигурационной абляции победила комбинация перекрёстной энтропии по меткам эмоций и KL-дистилляции от аудио к видео; подход не требует эмоциональных меток на инференсе.

Валидация --- RAVDESS, актёрно-разделённое разбиение 544/96/96; внешняя оценка двумя FER-классификаторами, отобранными скринингом шести кандидатов; критерий межклассификаторного согласия защищает от артефактов выбора эвалюатора.

\textbf{Результаты.} Wav2Lip: $\Delta\text{F1}=+10{,}7$~\% на тесте (0,611$\to$0,718), LSE-C незначимо ($p=0{,}61$); бутстрап-95~\% CI на эвалюаторе Rajaram1996 не включает ноль; согласие двух внешних подтверждается для всех 8 эмоциональных конфигов. SadTalker: $\Delta\text{F1}=+14{,}3$~\% на коэффициент-уровне ($p=0{,}0095$); на rendered-уровне не подтверждён. Вклад ограничен двумя архитектурами, четырьмя эмоциями и студийным RAVDESS.

\textbf{Ключевые слова:} генерация говорящих лиц, кросс-модальная функция потерь, эмоциональная согласованность, дообучение, KL-дистилляция, Wav2Lip, SadTalker, RAVDESS.

\newpage
% \renewcommand{\contentsname}{Оглавление}
\tableofcontents  % содержание
\suppressfloats
\newpage

\section{Введение}

\subsection{Актуальность исследования}

Технологии генерации говорящих лиц составляют один из наиболее быстрорастущих сегментов цифрового медиапроизводства и применяются в цифровых ассистентах, дубляже и локализации видеоконтента, образовательных платформах и виртуальных тьюторах, телемедицинских приложениях, игровой индустрии и инструментах AR/VR. По прогнозу Grand View Research, мировой рынок цифровых аватаров возрастёт с 18 миллиардов долларов США в 2023 году до 270 миллиардов долларов США к 2030 году с совокупным годовым темпом роста около 47~\% \cite{gvr2024}. Указанная динамика обусловливает повышение требований к естественности эмоционального поведения генерируемых аватаров: с расширением охвата применений возрастает и доля сценариев, в которых пользователи взаимодействуют с системой в эмоционально насыщенном контексте --- от обучения детей с особенностями развития до консультаций по психологическому здоровью и сопровождения пациентов с когнитивными нарушениями.

Современные архитектуры аудио-управляемого синтеза говорящих лиц достигли высоких показателей точности артикуляции. Модель Wav2Lip \cite{prajwal2020wav2lip} обеспечивает точность синхронизации губ свыше 97~\% по метрикам LSE-C/LSE-D на стандартных тестовых корпусах при работе с произвольным лицом. Модель SadTalker \cite{zhang2023sadtalker} реализует генерацию трёхмерной геометрии лица и движений головы с применением промежуточной 3DMM-репрезентации, что обеспечивает более естественные мимические переходы по сравнению с image-space подходами. Однако перечисленные модели ориентированы исключительно на синхронизацию артикуляции и не предусматривают явного контроля соответствия эмоционального содержания речи и мимики. Вследствие этого при подаче эмоционально насыщенного аудиосигнала генерируемое видео сохраняет нейтральное либо слабовыраженное мимическое выражение, что приводит к рассогласованию между смысловой и эмоциональной составляющей мультимодального вывода и снижает воспринимаемый реализм генерируемого контента.

Перцепционные исследования последовательно демонстрируют, что эмоциональная согласованность мультимодального вывода критически влияет на воспринимаемый реализм аватара и точность интерпретации эмоций пользователем. Работа \cite{suma2023} фиксирует снижение точности распознавания эмоций при рассогласовании голосового и мимического каналов на 14--22~\% относительно согласованной подачи; в \cite{zhou2025} эффект подтверждён на расширенной выборке с участием 240 респондентов в задаче оценки реализма синтезированных видео. Классическое наблюдение мультимодальной интеграции в восприятии речи \cite{mcgurk1976} подтверждает невозможность сведения восприятия мультимодального стимула к независимой обработке отдельных каналов; следовательно, генеративные системы, оптимизированные исключительно по одному каналу (артикуляции), систематически уступают системам, обеспечивающим согласованность по всем каналам, в субъективных оценках реализма. Дополнительно следует отметить эффект <<зловещей долины>> (uncanny valley), описанный в работе \cite{mori2012}: системы, демонстрирующие промежуточную, но рассогласованную натуралистичность, могут вызывать выраженный негативный отклик у пользователя, что в практическом контексте делает рассогласование эмоциональных сигналов прямой причиной отказа от использования системы.

В период с 2022 по 2026 год опубликован ряд работ, в которых эмоционально-обусловленное управление является существенным компонентом метода. EmoTalk \cite{peng2023emotalk} выполняет дисбэндлинг эмоциональных и лингвистических признаков для трёхмерной анимации лица; EmoTalker \cite{shen2024emotalker} применяет кросс-модальный трансформер с эмоциональным пулингом признаков; EmotiveTalk \cite{wang2025emotivetalk} реализует диффузионный подход с явной декомпозицией аудиоканала; EAT \cite{wang2023eat} использует эмоциональные адаптеры на основе MEAD; EmoCAST \cite{jiang2025emocast} предлагает диффузионный фреймворк с управлением через текстовое описание эмоции; AUHead \cite{lyu2026auhead} применяет промежуточное Action Unit представление; EmoDiffTalk \cite{liu2026emodifftalk} комбинирует AU-обусловливание с трёхмерный гауссов сплаттинг. Перечисленные подходы либо требуют замены опорной архитектуры (диффузионная опорная модель или специализированный энкодер-декодер), либо предполагают явную разметку эмоций или AU на этапе инференса. Применимость указанных подходов к задаче постфактумной адаптации существующих развёрнутых систем синтеза говорящих лиц ограничена: переход на специализированные эмоциональные модели требует существенных архитектурных и инфраструктурных изменений, переобучения сопутствующих компонентов пайплайна и пересборки сервисных интеграций.

Возможность повышения эмоциональной согласованности при сохранении уже развёрнутых опорных архитектур и без явных меток на этапе генерации в литературе системно не исследована. Сужение теоретического пробела до указанной постановки --- надстраиваемая кросс-модальная функция потерь, применимая в режиме файнтюнинга к двум репрезентативным архитектурно различным опорным моделям, без модификации архитектуры и без эмоциональных меток на инференсе --- определяет направление настоящего исследования. Дополнительная актуальность связана с методологическим аспектом оценки: в подавляющем большинстве работ по эмоциональному синтезу финальная оценка эмоциональной согласованности проводится теми же классификаторами эмоций, которые использованы в обучающей сигнальной цепи, что создаёт систематический риск кругового рассуждения \cite{jiang2024survey}. Настоящее исследование адресует указанный риск через предварительный скрининг внешних классификаторов и протокол межклассификаторное согласие, формулировка которого составляет самостоятельный методологический вклад.

\subsection{Объект и предмет исследования}

\textbf{Объект исследования} - модели аудио-управляемой генерации говорящих лиц Wav2Lip и SadTalker.

\textbf{Предмет исследования} - кросс-модальная функция потерь, связывающая эмоциональные представления речевого сигнала и видеоряда в едином векторном пространстве.

\subsection{Цель и задачи исследования}

\textbf{Цель исследования} - разработка и экспериментальная валидация кросс-модальной функции потерь, обеспечивающей статистически значимое улучшение F1-меры эмоциональной согласованности между речевым сигналом и мимикой при файнтюнинге опорных архитектур Wav2Lip и SadTalker на подмножестве из четырёх эмоций (\textit{happy}, \textit{sad}, \textit{angry}, \textit{disgust}) датасета RAVDESS без явной разметки эмоций на этапе генерации и без модификации архитектуры опорных моделей.

Для достижения поставленной цели решаются следующие задачи:
\begin{enumerate}
\item теоретический анализ существующих подходов к эмоционально-обусловленному синтезу говорящих лиц и формулировка теоретического пробела;
\item файнтюнинг аудио- и видеоэнкодеров эмоций на датасете RAVDESS как опорных моделей для извлечения эмоциональных представлений из соответствующих модальностей;
\item формализация кросс-модальной функции потерь и её абляционный анализ в пятиконфигурационном эксперименте на Wav2Lip с целью эмпирического выбора компонент (косинусное сближение, перекрёстная энтропия, KL-дистилляция и их сочетания);
\item файнтюнинг моделей Wav2Lip и SadTalker с применением победившей в абляции функции потерь и поиск весового коэффициента, обеспечивающего максимум F1 в пределах допустимой деградации реконструкции;
\item сравнительный анализ с базовыми моделями на тестовой выборке с проверкой статистической значимости (критерий Макнемара, парный $t$-критерий) и независимой внешней оценкой эмоциональной согласованности с использованием классификаторов, не задействованных в обучающей сигнальной цепи.
\end{enumerate}

\subsection{Гипотезы исследования}

\textbf{Гипотеза H1 (основная).} Включение предложенной кросс-модальной функции потерь в процедуру файнтюнинга Wav2Lip на четырёх эмоциях RAVDESS приводит к улучшению F1-меры эмоциональной согласованности не менее чем на 10~\% по сравнению с базовой моделью при условии, что абсолютное изменение метрики LSE-C не превышает 2~\% от значения базовой модели. \textit{Критерии подтверждения:} $\Delta\text{F1} \geq +10$~\% ($p < 0{,}05$, критерий Макнемара); $|\Delta\text{LSE-C}| \leq 2$~\% либо изменение LSE-C статистически незначимо по парному $t$-критерию ($p \geq 0{,}05$). \textit{Метрики снимаются на test-выборке RAVDESS из $n=96$ сэмплов; F1 оценивается дополнительно двумя независимыми внешними классификаторами выражений лица.}

\textbf{Гипотеза H2 (дополнительная).} Включение той же функции потерь в процедуру файнтюнинга SadTalker приводит к улучшению F1-меры эмоциональной согласованности на тестовой выборке не менее чем на 8~\% при сохранении ограничения на качество движений головы и общую визуальную целостность кадра, что демонстрирует методологическую воспроизводимость подхода на второй опорной архитектуре, отличающейся от Wav2Lip по типу представления (трёхмерная геометрия лица вместо двумерной маски). \textit{Критерии подтверждения:} $\Delta\text{F1} \geq +8$~\% ($p < 0{,}05$, критерий Макнемара) на тестовом подмножестве, оцениваемом независимыми внешними классификаторами по рендеренным видео.

\subsection{Методы исследования}

В работе применяются методы глубокого обучения (предобученные модели wav2vec~2.0 в варианте \texttt{superb/wav2vec2-base-superb-er} и TimeSformer в варианте \texttt{facebook/timesformer-base-finetuned-k400} для извлечения эмоциональных представлений из аудио- и видеомодальностей), а также методы кросс-модальной регуляризации, оптимальная конфигурация которой эмпирически устанавливается в пятиконфигурационной абляции. Для проверки статистической значимости применяется критерий Макнемара для категориальных метрик и парный $t$-критерий для метрик с количественной шкалой. Программная реализация выполнена на языке Python с применением библиотек PyTorch и Hugging Face Transformers; логирование экспериментов проведено в системе Weights \& Biases. Для снижения риска цикличности оценки применяется независимая внешняя оценка двумя классификаторами эмоций, не задействованными на этапах обучения и отбора моделей; оба внешних классификатора отобраны по результатам скрининга шести кандидатов на тестовой выборке RAVDESS.

\subsection{Научная новизна}

Научная новизна исследования формулируется в границах проведённого эксперимента и состоит в следующем. Впервые показано, что повышение эмоциональной согласованности генерации говорящих лиц на четырёх эмоциях RAVDESS достижимо посредством надстраиваемой функции потерь, применяемой в режиме файнтюнинга к опорной архитектуре Wav2Lip без модификации опорной архитектуры и без явных эмоциональных меток на этапе инференса; обоснованность выбора компонент функции потерь (победила комбинация перекрёстной энтропии и KL-дистилляции от аудиоэнкодера к видеоэнкодеру) подтверждена пятиконфигурационным абляционным анализом, в котором наивная косинусная формулировка показала отрицательную $\Delta\text{F1}$ относительно baseline. Эффект подтверждён двумя независимыми внешними классификаторами, не задействованными в обучающей сигнальной цепи и предварительно отобранными в процедуре скрининга. В отличие от EmoTalk, EmoTalker, EmotiveTalk и EAT, использующих специализированные архитектуры или диффузионные стеки и оценивающих результат на тех же моделях, что использованы в loss, предложенный подход совместим с уже развёрнутыми моделями и реализует протокол внешней оценки, систематически отсутствующий в перечисленных работах. Распространение подхода на SadTalker как надстройки над опорной моделью в текущем эксперименте не подтверждено; пределы переноса методологии на иные архитектуры обсуждаются в разделе перспектив.

\subsection{Теоретическая и практическая значимость}

\textbf{Теоретическая значимость} заключается в формализации задачи повышения эмоциональной согласованности при файнтюнинге как задачи кросс-модального контрастивного выравнивания на фиксированных энкодерах и в постановке протокола независимой внешней валидации, отделяющего вклад функции потерь от свойств обучающих метрик.

\textbf{Практическая значимость} определяется применимостью разработанной функции потерь к двум распространённым опорным архитектурам без переобучения с нуля, что снижает вычислительные затраты на адаптацию существующих систем и обеспечивает совместимость с уже развёрнутыми пайплайнами генерации говорящих лиц.

\subsection{Апробация результатов исследования}

Основные результаты диссертационного исследования представлены на 68-й Всероссийской научной конференции МФТИ (секция когнитивных технологий, 4 апреля 2026 г., Москва); работа принята для публикации в сборнике тезисов конференции (\url{https://conf.mipt.ru/conference/23659}). Программная реализация опубликована в открытом репозитории GitHub (\url{https://github.com/Katrin-Pochtar/The-Uncanny-Valley}); журналы экспериментов доступны в системе Weights \& Biases.

\subsection{Структура работы}

Диссертация состоит из введения, теоретического анализа, методологии исследования, главы результатов, перспектив и заключения, а также списка использованных источников и приложений. В теоретическом анализе систематизируются существующие подходы к синтезу говорящих лиц и формулируется теоретический пробел. В методологической главе описаны файнтюнинг аудио- и видеоэнкодеров, формализация и реализация кросс-модальной функции потерь, а также процедуры файнтюнинга Wav2Lip и SadTalker. В главе результатов представлены количественные результаты по гипотезам H1 и H2, абляционный анализ компонент функции потерь и независимая внешняя оценка. В разделе перспектив обсуждаются ограничения исследования и направления дальнейших работ. Заключение содержит обобщённые выводы по результатам проверки выдвинутых гипотез.

\section{Обозначения и сокращения}

\renewcommand{\arraystretch}{1.2}
\renewcommand{\tabcolsep}{0.3cm}
\begin{longtable}{|p{3.5cm}|p{12cm}|}
\hline
\textbf{Сокращение} & \textbf{Расшифровка} \\
\hline
\endfirsthead
\hline
\textbf{Сокращение} & \textbf{Расшифровка} \\
\hline
\endhead
ВКР & выпускная квалификационная работа \\
\hline
3DMM & трёхмерная морфируемая модель лица (3D Morphable Model) \\
\hline
AU & элементарная единица мимики (Action Unit) \\
\hline
CE & перекрёстная энтропия (Cross Entropy) \\
\hline
CNN & свёрточная нейронная сеть (Convolutional Neural Network) \\
\hline
FER & распознавание выражений лица (Facial Expression Recognition) \\
\hline
F1 & F1-мера, среднее гармоническое точности и полноты \\
\hline
GAN & генеративно-состязательная сеть (Generative Adversarial Network) \\
\hline
HuBERT & Hidden-Unit BERT - аудиоэнкодер с самообучением \\
\hline
KL & расхождение Кульбака - Лейблера (Kullback - Leibler divergence) \\
\hline
L1 & L1-норма ошибки реконструкции \\
\hline
LSE-C, LSE-D & Lip Sync Error - Confidence/Distance, метрики синхронизации губ \\
\hline
MEAD & Multi-view Emotional Audio-visual Dataset \\
\hline
MSP-DIM & модель размерной оценки эмоции (arousal/dominance/valence) \\
\hline
RAVDESS & Ryerson Audio-Visual Database of Emotional Speech and Song \\
\hline
SupCon & Supervised Contrastive Loss (Khosla et al., 2020) \\
\hline
SyncNet & модель оценки качества синхронизации губ \\
\hline
TimeSformer & пространственно-временной трансформер для видео \\
\hline
ViT & Vision Transformer \\
\hline
VAD & размерное представление эмоции: валентность/возбуждение/доминирование \\
\hline
VAE / VQ-VAE & вариационный автокодировщик / его векторно-квантованный вариант \\
\hline
Wav2Lip & модель аудио-управляемой синхронизации губ \\
\hline
Wav2Vec2 & самообучающаяся аудио-модель \\
\hline
W\&B & сервис логирования экспериментов Weights \& Biases \\
\hline
$\Delta$F1 & изменение F1-меры относительно базовой модели \\
\hline
$\Delta$LSE-C & изменение метрики LSE-C относительно базовой модели \\
\hline
H1, H2 & основная и дополнительная гипотезы исследования \\
\hline
\end{longtable}

\section{Теоретический анализ}

\subsection{Логика построения обзора литературы}

Настоящий раздел реализует системный обзор научной литературы по тематике аудио-управляемого синтеза говорящих лиц с акцентом на проблему эмоциональной согласованности. Структура обзора подчинена логике последовательного сужения от широкого контекста к узкой исследовательской постановке. На первом этапе рассматриваются психологические основы классификации эмоций, определяющие выбор представления эмоционального пространства. На втором этапе анализируются основные архитектурные классы существующих генеративных моделей и характер используемых ими функций потерь. На третьем этапе систематизируются подходы к интеграции эмоционального компонента в задачу синтеза говорящих лиц с разделением на категории по типу управляющего сигнала и степени модификации базовой архитектуры. На четвёртом этапе анализируются предобученные энкодеры эмоциональных представлений для аудио- и видеомодальностей. На пятом этапе рассматриваются датасеты, используемые в задачах распознавания эмоций и эмоционально-обусловленного синтеза. На шестом этапе рассматриваются методы кросс-модального контрастивного обучения, составляющие теоретическую основу предлагаемой функции потерь. На седьмом этапе анализируются метрики оценки качества генерации и эмоциональной согласованности с особым вниманием к методологическому риску цикличности оценки. На заключительном этапе проводится систематизация выявленных подходов в форме сравнительной таблицы и формулируется теоретический пробел, обосновывающий направление настоящего исследования.

\subsection{Психологические основы классификации эмоций}

Выбор представления эмоционального пространства в задачах распознавания и синтеза эмоционально-окрашенного контента опирается на классические психологические модели эмоций.

\textbf{Категориальная модель базовых эмоций.} Согласно теории, формализованной в работах \cite{ekman1992,ekman1978facs}, существует ограниченный набор универсальных базовых эмоций, имеющих кросс-культурное проявление и характерные мимические корреляты. В классической формулировке выделяются шесть базовых эмоций (радость, печаль, гнев, страх, отвращение, удивление); расширенный вариант включает дополнительно нейтральное и презрительное состояния. Категориальная модель составляет основу большинства корпусов с эмоциональной разметкой (RAVDESS, CREMA-D, AffectNet) и применяется в качестве стандарта для оценки моделей распознавания и генерации эмоций.

\textbf{Размерная модель.} Альтернативный подход, восходящий к работам \cite{russell1980} и развитый в модели VAD (valence-arousal-dominance), представляет эмоциональное состояние не дискретной меткой, а точкой в непрерывном многомерном пространстве. Валентность характеризует приятность переживания, возбуждение - интенсивность активации, доминирование - степень контроля. Размерная модель обеспечивает более тонкую градацию состояний и применяется в корпусах MSP-Podcast и SEMAINE; в задачах распознавания речевых эмоций реализована в моделях \cite{wagner2023}. Преимуществом размерной модели является непрерывная природа представления, недостатком - менее интуитивная интерпретация и сложность построения корпусов с надёжной разметкой.

\textbf{Колесо эмоций Плутчика.} Модель \cite{plutchik1980} вводит восемь базовых эмоций, попарно противоположных, и допускает построение производных смешанных состояний. Указанная модель используется преимущественно в задачах анализа текста и менее распространена в аудиовизуальном анализе.

\textbf{Выбор представления для настоящего исследования.} В работе используется категориальная модель базовых эмоций в варианте, ограниченном четырьмя классами датасета RAVDESS: радость (happy), печаль (sad), гнев (angry) и отвращение (disgust). Сужение исходного восьмиклассового пространства RAVDESS до четырёх категорий мотивировано двумя соображениями. Во-первых, состояния <<нейтральное>> (neutral) и <<спокойное>> (calm) обладают слабой эмоциональной выраженностью и низкой различимостью при автоматической оценке мимики, что делает их включение методологически проблемным в задаче, где целевой эффект - именно повышение эмоционально окрашенной выразительности. Во-вторых, классы <<страх>> (fearful) и <<удивление>> (surprised) в студийной актёрской записи RAVDESS систематически конфузируются с радостью на уровне внешних предобученных классификаторов FER-производного происхождения, использование которых составляет ядро протокола внешней оценки в настоящей работе; включение указанных классов привело бы к интерференции между источниками шума оценивания и истинным эффектом функции потерь. Используемый набор из четырёх эмоций согласуется с распространённой практикой работ по эмоциональному синтезу \cite{ji2022eamm,wang2023eat} и обеспечивает интерпретируемое сопоставление матриц ошибок между внутренними и внешними классификаторами. Ограничения, связанные с указанным сужением эмоционального пространства, явно отражены в разделе, посвящённом границам исследования.

\subsection{Архитектурные классы моделей генерации говорящих лиц}

Современные модели аудио-управляемого синтеза говорящих лиц могут быть разделены на четыре архитектурных класса по способу представления и обновления лицевой геометрии: двумерные модели (image-space), трёхмерные модели на основе морфируемых представлений лица (3DMM-ориентированный), модели на основе нейронных радиационных полей и трёхмерного гауссова сплаттинга (NeRF / трёхмерный гауссов сплаттинг) и диффузионные модели.

\textbf{Двумерные image-space модели} оперируют непосредственно в пространстве пикселей входного кадра. Архитектура Wav2Lip \cite{prajwal2020wav2lip} основана на парадигме <<кодировщик - декодер>> с дополнительным дискриминатором синхронизации губ (SyncNet), что обеспечивает точность синхронизации артикуляции свыше 97~\% на стандартных тестовых корпусах при работе с произвольным лицом. Модель оптимизирует функцию реконструкции L1 в нижней половине лица и совместно обучается с предобученным lipsync-экспертом, замороженным в процессе тренировки. Дальнейшие модификации, включая Wav2Lip-HQ \cite{mallikarjuna2025wav2liphq} и VividWav2Lip \cite{liu2024vividwav2lip}, направлены на повышение визуального качества посредством сегментации лица и применения сверхразрешающих сетей; явных механизмов управления эмоциональной выразительностью в указанных модификациях не вводится. К данному классу также относится модель MakeItTalk \cite{zhou2020makeittalk}, реализующая управление позой головы через предсказание ключевых точек лица в качестве промежуточного представления.

\textbf{Трёхмерные модели на основе морфируемой геометрии} разделяют процесс синтеза на два этапа: на первом из аудиосигнала предсказываются параметры трёхмерной геометрии лица (выражение, поза, освещение), на втором рендерятся видеокадры. Модель SadTalker \cite{zhang2023sadtalker} использует двухветвевой подход: блок PoseVAE предсказывает движение головы, блок ExpNet - параметры выражений, после чего применяется промежуточная 3DMM-репрезентация. Подобная декомпозиция обеспечивает более естественные движения головы по сравнению с image-space моделями, однако эмоциональное содержание контролируется только опосредованно через параметры выражений, без явного связывания с эмоциональной семантикой аудиосигнала. Сходный двухэтапный подход применён в моделях Audio2Head \cite{wang2021audio2head}, EAMM \cite{ji2022eamm} и DiffPoseTalk \cite{sun2024diffposetalk}. Общим ограничением класса является зависимость качества выходного видео от точности промежуточной геометрической репрезентации.

\textbf{Модели на основе нейронных радиационных полей} реализуют синтез фотореалистичного видео через объёмное представление сцены, обусловленное аудиопризнаками. Семейство AD-NeRF \cite{guo2021adnerf}, RAD-NeRF \cite{tang2022radnerf} и GeneFace \cite{ye2023geneface} требует индивидуального обучения NeRF на видеоматериале конкретного субъекта, что ограничивает обобщение на произвольное лицо, но обеспечивает высокую визуальную точность для целевого субъекта. Подходы на основе трёхмерный гауссов сплаттинг (InsTaG \cite{li2025instag}, EmoDiffTalk \cite{liu2026emodifftalk}) сокращают время обучения и улучшают параметрическую управляемость, в том числе через интеграцию управления Action Units. Эмоциональный контроль в данном классе моделей реализуется преимущественно через явные текстовые или AU-промпты на этапе генерации, что требует наличия дополнительной разметки и не соответствует постановке задачи извлечения эмоции непосредственно из аудиосигнала.

\textbf{Диффузионные модели} реализуют синтез как последовательный процесс денойзинга, обусловленный аудиопризнаками. Модель DiffTalk \cite{shen2023difftalk} применяет латентное диффузионное моделирование для портретной анимации; EMO \cite{tian2024emo} комбинирует диффузионный процесс с временной свёрткой для генерации длительных эмоционально насыщенных последовательностей; EmotiveTalk \cite{wang2025emotivetalk} вводит декуплирование аудиопризнаков для раздельного контроля артикуляции и эмоций. Диффузионные подходы достигают высокого качества генерации, однако требуют существенных вычислительных ресурсов и, как правило, замены всего генеративного стека, что ограничивает их применимость к задаче файнтюнинга уже развёрнутых моделей.

Систематизация показывает, что несмотря на архитектурное разнообразие, перечисленные классы моделей либо не предусматривают явного управления эмоциональной выразительностью (классы 1 и 2), либо требуют замены архитектуры или явных эмоциональных промптов на этапе инференса (классы 3 и 4). Опорные архитектуры Wav2Lip и SadTalker, выбранные в качестве объектов настоящего исследования, относятся к первым двум классам и представляют собой широко применяемые в практике решения, для которых задача надстройки эмоционального контроля без модификации архитектуры представляет самостоятельный исследовательский интерес.

\subsection{Эмоциональное измерение в синтезе говорящих лиц}

Подходы к интеграции эмоционального компонента в задачу синтеза говорящих лиц могут быть классифицированы по типу управляющего сигнала и по степени модификации базовой архитектуры. Указанная классификация позволяет позиционировать настоящее исследование относительно существующих работ.

\textbf{Управление через явные категориальные метки} реализовано в семействе моделей, обученных на корпусах с эмоциональной разметкой. Датасет MEAD \cite{wang2020mead} (Multi-view Emotional Audio-visual Dataset) содержит видеозаписи 60 актёров с восемью базовыми эмоциями и тремя уровнями интенсивности; на его основе разработаны специализированные эмоциональные модели, включая EAMM \cite{ji2022eamm} и EAT \cite{wang2023eat}. Модель EAT (Efficient Emotional Adaptation for Audio-Driven Talking-Head Generation) реализует адаптацию через дополнительные эмоциональные модули, требующие явной метки эмоции на этапе инференса. Подобный подход обеспечивает прямой контроль выразительности, но требует дополнительной разметки в применении и не масштабируется на сценарии, где эмоция должна выводиться непосредственно из аудиосигнала.

\textbf{Управление через специализированную архитектуру} реализовано в моделях, проектируемых с эмоциональным компонентом как часть целевой функции обучения. EmoTalk \cite{peng2023emotalk} выполняет разделение эмоциональных и лингвистических признаков для трёхмерной анимации лица посредством специальных энкодеров и обучается с нуля. EmoCAST \cite{jiang2025emocast} применяет диффузионный фреймворк с управлением через текстовые описания эмоций. EMOCA \cite{danecek2022emoca} решает обратную задачу - реконструкцию трёхмерной геометрии лица из видео с учётом эмоциональной информации - и не предназначена для аудио-управляемого синтеза. Указанные подходы либо требуют полной переработки архитектуры, либо относятся к иной постановке задачи.

\textbf{Управление через Action Units} опирается на систему кодирования лицевых действий (Facial Action Coding System, FACS), формализованную в работе \cite{ekman1978facs}. Action Units представляют собой элементарные мимические движения, имеющие нейроанатомическую интерпретацию и обеспечивающие компактное структурированное описание выражения лица. Модели AUHead \cite{lyu2026auhead}, MF-ETalk \cite{wang2025emotivetalk} и EmoDiffTalk \cite{liu2026emodifftalk} используют AU как промежуточное представление между аудиосигналом и видеовыходом. AU-предсказывающие модели обеспечивают интерпретируемость и анатомическую корректность, однако требуют вычислительно дорогих этапов промежуточной разметки и систем предсказания AU из аудио.

\textbf{Управление через кросс-модальное выравнивание без явных меток на инференсе} представлено меньшим числом работ. EmoTalker \cite{shen2024emotalker} реализует кросс-модальный трансформер, обученный с нуля, который извлекает эмоциональную информацию непосредственно из аудио и применяет её при генерации без необходимости в явной метке. CEM-Net \cite{wu2025cemnet} решает связанную задачу разрешения конфликта между референсным изображением и аудиосигналом посредством сети памяти. Подобные подходы концептуально близки к настоящему исследованию, однако предполагают обучение специализированной архитектуры с нуля, а не файнтюнинг существующей опорной модели.

\textbf{Полемика в литературе.} В исследовательском сообществе наблюдаются различные позиции относительно компромисса между интерпретируемостью и обобщающей способностью. Сторонники AU-управления \cite{lyu2026auhead,liu2026emodifftalk} аргументируют необходимость явного структурированного представления для обеспечения контроля и анатомической корректности, опираясь на нейропсихологическую обоснованность системы FACS. Сторонники end-to-end кросс-модальных подходов \cite{shen2024emotalker,wang2025emotivetalk} указывают на ограниченность дискретных категориальных или AU-представлений для передачи тонких эмоциональных оттенков и отстаивают применение непрерывных аудио-производных представлений. Перцепционные исследования \cite{suma2023,zhou2025} подтверждают значимость согласованности голосового и мимического каналов независимо от конкретного механизма управления, опираясь на классическое наблюдение мультимодальной интеграции в восприятии речи \cite{mcgurk1976}.

Систематизация выявляет отсутствие в литературе работ, в которых эмоциональная согласованность достигается посредством надстройки над уже существующими опорными архитектурами без их замены и без явной разметки эмоций на этапе инференса. Указанная ниша определяет постановку настоящего исследования.

\subsection{Аудио- и видеоэнкодеры эмоциональных представлений}

Извлечение эмоциональных представлений из соответствующих модальностей реализуется специализированными предобученными моделями, обучение которых на корпусах с эмоциональной разметкой обеспечивает высокую разделяющую способность в пространстве эмбеддингов.

\textbf{Аудиоэнкодеры.} Архитектуры на основе самообучения с маскировкой (self-supervised learning) показали высокую эффективность в задачах распознавания речевых эмоций. Модель wav2vec~2.0 \cite{baevski2020wav2vec} обучается на сырых аудиоволнах через контрастивную задачу предсказания квантованных представлений; HuBERT \cite{hsu2021hubert} заменяет контрастивную задачу на маскированное предсказание кластерных меток, полученных кластеризацией k-means на признаках MFCC. После файнтюнинга на эмоциональных корпусах (RAVDESS, IEMOCAP) указанные модели достигают F1-меры 0,75-0,90 на классификации базовых эмоций \cite{pepino2021}. Расширение AV-HuBERT \cite{shi2022avhubert} обобщает HuBERT на аудиовизуальную модальность через совместное маскированное обучение и применяется в задачах распознавания речи в шумных условиях.

\textbf{Видеоэнкодеры.} Архитектура Vision Transformer \cite{dosovitskiy2021vit} обобщает механизм самовнимания на изображения через разбиение на патчи и линейное проецирование. TimeSformer \cite{bertasius2021timesformer} адаптирует ViT к видео через факторизацию пространственного и временного внимания, что снижает вычислительную сложность и обеспечивает применимость на длинных последовательностях кадров. Для распознавания выражений лица применяются специализированные ViT, предобученные на корпусах AffectNet \cite{mollahosseini2017affectnet} и FER+ \cite{barsoum2016ferplus}. Альтернативный класс представлен моделями MARLIN \cite{cai2023marlin} и Video Masked Autoencoders \cite{tong2022videomae}, реализующими маскированное автокодирование для видеомодальности. Для задач, требующих интерпретируемых выражений, применяются также модели предсказания AU, обученные на корпусах с FACS-разметкой.

\textbf{Кросс-модальные представления.} CLIP \cite{radford2021clip} продемонстрировал эффективность контрастивного обучения на парах <<изображение - текст>> для построения общего семантического пространства. ALIGN \cite{jia2021align} масштабирует подход на корпусе из 1,8 миллиарда пар. CAV-MAE \cite{gong2023cavmae} адаптирует масочное автокодирование к аудиовизуальной модальности и реализует совместное обучение через комбинацию реконструкции и контрастивного выравнивания. Применение принципа CLIP к паре <<эмоциональное аудио - эмоциональное видео>> при синтезе говорящих лиц в литературе систематически не исследовано; имеющиеся работы в области аудио-визуального согласования (например, \cite{owens2018}) сфокусированы на задачах распознавания речи и идентификации говорящего, а не на эмоциональной семантике в задаче генерации.

\textbf{Выбор энкодеров для настоящего исследования.} В качестве отправной точки рассмотрены три семейства предобученных аудиомоделей (Wav2Vec2 \texttt{base}, Wav2Vec2 \texttt{large}, HuBERT \texttt{large}, а также SUPERB-варианты, дообученные на корпусе IEMOCAP) и одно семейство видеомоделей (TimeSformer, варианты с числом кадров 4/8/12/16 и различными темпами обучения). По итогам сравнительной оценки на валидационной выборке RAVDESS (4-классовая постановка, разбиение по актёрам) для дальнейшей работы отобраны: \texttt{superb/wav2vec2-base-superb-er} в качестве аудиоэнкодера (валидационная макро-F1 0,826) и \texttt{facebook/timesformer-base-finetuned-k400} с 8-кадровой выборкой в качестве видеоэнкодера (валидационная макро-F1 0,790). На независимой тестовой выборке указанные модели достигают макро-F1 0,739 (аудио) и 0,876 (видео), что фиксирует <<потолок>> эмоциональной разделимости, доступной любой модели генерации, оцениваемой через данные энкодеры. Подробности гиперпараметрического перебора и полная таблица конфигураций приведены в разделе результатов и в ноутбуке [02\_train\_emotion\_encoders.ipynb](02\_train\_emotion\_encoders.ipynb). Выходы классификационных голов указанных моделей (4-мерные логиты эмоций) используются как целевые сигналы в формулируемой далее кросс-модальной функции потерь. Для оценки на этапе финального тестирования применяются независимые внешние классификаторы выражений лица, не задействованные ни на одном этапе обучения и отбора моделей; обоснование выбора и процедура скрининга приведены в подразделе, посвящённом метрикам.

\subsection{Датасеты для эмоционально-обусловленного синтеза}

Качество эмоционально-обусловленных моделей генерации существенно зависит от характеристик обучающих корпусов: количества эмоциональных категорий, числа и разнообразия субъектов, условий записи и наличия парной аудиовизуальной разметки. В литературе сложился ограниченный круг датасетов, регулярно применяемых в задачах распознавания и синтеза эмоций.

\textbf{RAVDESS (Ryerson Audio-Visual Database of Emotional Speech and Song)} \cite{livingstone2018ravdess} содержит 2880 аудиовизуальных записей 24 актёров (12 мужчин и 12 женщин), произносящих две стандартные фразы в восьми эмоциональных состояниях (нейтральное, спокойное, радость, печаль, гнев, страх, отвращение, удивление) с двумя уровнями интенсивности. Запись осуществлялась в студийных условиях с фиксированной позой головы и контролируемым освещением, что обеспечивает стандартизацию, но ограничивает внешнюю валидность относительно условий в естественных условиях. RAVDESS является наиболее широко используемым сбалансированным корпусом для аудиовизуального распознавания эмоций с парной разметкой речи и видео и применяется в качестве основного датасета настоящего исследования.

\textbf{MEAD (Multi-view Emotional Audio-visual Dataset)} \cite{wang2020mead} содержит видеозаписи 60 актёров с восемью базовыми эмоциями и тремя уровнями интенсивности, снятые с семи ракурсов. MEAD превосходит RAVDESS по объёму и разнообразию субъектов и применяется в специализированных эмоциональных моделях \cite{ji2022eamm,wang2023eat}. Использование MEAD в настоящем исследовании выходит за пределы экспериментально подтверждённого охвата и составляет одно из направлений дальнейших работ.

\textbf{CREMA-D (Crowd-sourced Emotional Multimodal Actors Dataset)} \cite{cao2014cremad} содержит 7442 записи 91 актёра разнородного этнокультурного происхождения с шестью базовыми эмоциями и четырьмя уровнями интенсивности. По сравнению с RAVDESS обеспечивает большее разнообразие субъектов; внешняя валидация на CREMA-D также рассматривается как направление дальнейших исследований.

\textbf{IEMOCAP (Interactive Emotional Dyadic Motion Capture)} \cite{busso2008iemocap} содержит около 12 часов аудиовизуальных диалогов с двойной разметкой (категориальной и размерной), полученной от нескольких аннотаторов. IEMOCAP применяется преимущественно в задачах распознавания речевых эмоций и менее подходит для задач синтеза говорящих лиц вследствие диалогового характера записи и отсутствия фиксированных стандартизованных фраз.

\textbf{AffectNet} \cite{mollahosseini2017affectnet} и \textbf{FER+} \cite{barsoum2016ferplus} являются крупномасштабными корпусами изображений лиц с эмоциональной разметкой (около 1 миллиона и 35 тысяч изображений соответственно). Указанные корпуса используются для предобучения видеоэнкодеров и независимых классификаторов выражений лица; в настоящем исследовании внешний классификатор основан на FER-производном корпусе, что обеспечивает независимость от RAVDESS.

\textbf{Обоснование выбора RAVDESS.} Выбор RAVDESS в качестве основного датасета мотивирован следующим: (1) сбалансированное представление восьми категорий эмоций при одинаковом числе записей на категорию обеспечивает корректное обучение и оценку без необходимости weighted sampling; (2) парная аудиовизуальная разметка с фиксированными фразами устраняет конфаундинг между лингвистическим содержанием и эмоциональной окраской; (3) контролируемые условия записи минимизируют влияние освещения и позы головы на извлекаемые эмоциональные представления; (4) корпус является открыто доступным и широко используемым в литературе, что обеспечивает воспроизводимость и сопоставимость результатов. Ограничения, связанные со студийным характером записи и относительно небольшим числом субъектов, явно отражены в разделе, посвящённом границам исследования.

\subsection{Контрастивное обучение и кросс-модальное выравнивание}

Теоретическую основу предлагаемой функции потерь составляют методы контрастивного обучения, цель которых - построение представлений, в которых семантически близкие объекты сближаются, а семантически различные - раздвигаются в общем векторном пространстве.

\textbf{Фундаментальные формулировки.} Целевая функция InfoNCE \cite{vandenoord2018infonce} обеспечивает построение representation через максимизацию нижней границы взаимной информации между связанными представлениями посредством классификации позитивной пары среди множества негативных. Метод SimCLR \cite{chen2020simclr} продемонстрировал применимость подхода к самообучению на изображениях, используя augmentations как источник позитивных пар. Метод MoCo \cite{he2020moco} вводит momentum encoder и dynamic queue для увеличения числа доступных негативных примеров, что особенно существенно при ограниченном размере батча.

\textbf{Supervised contrastive learning.} Метод SupCon \cite{khosla2020supcon} расширяет контрастивную постановку на supervised режим: позитивами для якоря $i$ являются все объекты с той же меткой класса, а не только augmentations того же объекта. Подобная формулировка обеспечивает построение более структурированного embedding-пространства и в ряде задач классификации показала превосходство над cross-entropy. Применение SupCon к кросс-модальной постановке (<<аудио - видео>>) с использованием эмоциональных меток как источника позитивных пар реализуется в настоящем исследовании в качестве основной формулировки функции потерь и сопоставляется в абляции с альтернативными вариантами.

\textbf{Кросс-модальное выравнивание.} CLIP \cite{radford2021clip} обучает совместное представление изображений и текста через симметричную InfoNCE-постановку с использованием в качестве позитивных пар связанных по природе объектов и в качестве негативных - несвязанных пар в батче. Применительно к аудио-видео модальности AV-HuBERT \cite{shi2022avhubert} и CAV-MAE \cite{gong2023cavmae} реализуют выравнивание через совместное маскированное обучение, однако не специализированы на эмоциональной семантике; используемые ими позитивные пары формируются по принципу синхронности во времени, а не по эмоциональной общности.

\textbf{Риск коллапса проекций.} При попытке наивного применения косинусного сближения к парам <<аудио - видео>> одного и того же сэмпла без использования негативных примеров возникает риск тривиального коллапса: проекционные слои сходятся к константному отображению, обеспечивающему максимум косинусного сходства без передачи семантической информации \cite{chen2021simsiam}. Механизмы предотвращения коллапса в литературе по самообучению включают применение остановленного градиента (stop-gradient), использование предсказателя на одной из ветвей (BYOL \cite{grill2020byol}) и явное введение негативных пар. В настоящем исследовании предотвращение коллапса достигается посредством перехода от наивной косинусной формулировки к контрастивной с явными негативными парами, формируемыми на основе несовпадения эмоциональных меток. Указанный методологический выбор подтверждён эмпирически в рамках абляционного анализа.

\subsection{Метрики оценки качества генерации и эмоциональной согласованности}

Оценка моделей синтеза говорящих лиц традиционно опирается на несколько групп метрик, характеризующих различные аспекты качества генерации.

\textbf{Метрики синхронизации губ.} Стандарт де-факто составляют метрики LSE-C (Lip-Sync Error - Confidence) и LSE-D (Lip-Sync Error - Distance), вычисляемые с применением модели SyncNet \cite{chung2016syncnet}. SyncNet представляет собой двухветвевую сеть, обученную на аудиовизуальных парах для предсказания временного смещения; LSE-C измеряет уверенность синхронизации, LSE-D - нормированное расстояние в общем embedding-пространстве. Метрики SyncNet используются в качестве первичного бенчмарка в работах \cite{prajwal2020wav2lip,zhang2023sadtalker,ye2023geneface} и составляют ограничение на допустимое ухудшение качества артикуляции, формализованное в гипотезах H1 и H2 настоящего исследования.

\textbf{Метрики визуального качества.} Frechét Inception Distance \cite{heusel2017fid} оценивает расстояние между распределениями признаков сгенерированных и эталонных изображений в пространстве предобученной InceptionV3. Frechét Video Distance \cite{unterthiner2018fvd} обобщает FID на видео через применение Inflated 3D ConvNet. LPIPS \cite{zhang2018lpips} оценивает перцепционное сходство в пространстве признаков предобученной нейронной сети с обучением калибровочных весов на парных данных с человеческими оценками. Указанные метрики характеризуют общее визуальное качество, но не специфичны к эмоциональной выразительности и не позволяют оценить корректность эмоциональной семантики сгенерированного видео.

\textbf{Метрики эмоциональной согласованности.} Стандартный подход состоит в применении предобученного классификатора эмоций к сгенерированному видео и сопоставлении предсказанной метки с целевой эмоцией; в качестве агрегирующей метрики используется F1-мера или accuracy \cite{ji2022eamm,wang2023eat}. В отличие от метрик синхронизации губ, для оценки эмоциональной согласованности в литературе не сложилось универсального бенчмарка, и конкретный выбор классификатора (его обучающие данные, архитектура, версия) существенно влияет на результат, что затрудняет прямое сравнение между работами.

\textbf{Проблема цикличности оценки.} Если эмоциональный классификатор, используемый в качестве компонента функции потерь при обучении (например, для извлечения целевых эмоциональных представлений), одновременно применяется и для финальной оценки эмоциональной согласованности сгенерированного видео, возникает риск кругового рассуждения: модель оптимизируется под распределение, которое тот же классификатор успешно распознаёт, что не гарантирует переноса на независимое восприятие. Указанное методологическое ограничение систематически отмечается в обзорной литературе \cite{jiang2024survey}, однако в практике многих работ по эмоциональному синтезу мер по его явному устранению не предусмотрено. В настоящем исследовании в качестве методологического ответа на данный риск проведён скрининг шести предобученных внешних классификаторов выражений лица (\texttt{dima806/facial\_emotions\_image\_detection}, \texttt{trpakov/vit-face-expression}, \texttt{motheecreator/vit-Facial-Expression-Recognition}, \texttt{Rajaram1996/FacialEmoRecog}, \texttt{RickyIG/emotion\_face\_image\_classification}, \texttt{prithivMLmods/Facial-Emotion-Detection-SigLIP2}); последний исключён из-за несовместимости меток (отсутствие класса \textit{disgust}). По результатам скрининга на тестовой выборке RAVDESS (см. раздел результатов) для финальной оценки отобраны два классификатора с максимальной макро-F1 на четырёх целевых эмоциях - \texttt{motheecreator/vit-Facial-Expression-Recognition} и \texttt{Rajaram1996/FacialEmoRecog}; оба обучены на FER-производных корпусах, не пересекающихся с RAVDESS, что обеспечивает независимость оценочного сигнала. Перенос направления эффекта между двумя различающимися по архитектуре и обучающему корпусу внешними классификаторами фиксируется в виде отдельной диагностической метрики (межклассификаторное согласие).

\textbf{Перцепционные метрики.} Слепое A/B тестирование и оценка по шкале Ликерта \cite{likert1932} остаются золотым стандартом для оценки субъективного качества генерации \cite{zhou2025}. Объёмная перцепционная оценка с привлечением группы респондентов выходит за пределы настоящего исследования и составляет одно из направлений дальнейших работ, обозначенных в разделе перспектив.

\subsection{Систематизация и формулировка теоретического пробела}

Проведённый обзор позволяет систематизировать существующие подходы к эмоционально-обусловленной генерации говорящих лиц по нескольким критериям, существенным для настоящего исследования: требуется ли замена базовой архитектуры (опорной модели); требуется ли явная метка эмоции на этапе инференса; специализирован ли подход на эмоциональной семантике; реализована ли независимая внешняя валидация. Результаты систематизации представлены в таблице~\ref{tab:approaches_comparison}.

\renewcommand{\arraystretch}{1.3}
\renewcommand{\tabcolsep}{0.15cm}
\begin{table}[H]
\caption{Сравнение подходов к эмоционально-обусловленной генерации говорящих лиц}
\label{tab:approaches_comparison}
\centering
\small
\begin{tabular}{|p{4.5cm}|c|c|c|c|}
\hline
\textbf{Подход} & \textbf{Замена опорной модели} & \textbf{Явные метки на инференсе} & \textbf{Эмоциональная специализация} & \textbf{Внешняя валидация} \\
\hline
Wav2Lip \cite{prajwal2020wav2lip} & - & - & нет & - \\
\hline
SadTalker \cite{zhang2023sadtalker} & - & - & нет & - \\
\hline
EmoTalk \cite{peng2023emotalk} & да & да & да & нет \\
\hline
EAT \cite{wang2023eat} & да & да & да & нет \\
\hline
EmoTalker \cite{shen2024emotalker} & да & нет & да & нет \\
\hline
EmotiveTalk \cite{wang2025emotivetalk} & да & нет & да & нет \\
\hline
EmoCAST \cite{jiang2025emocast} & да & да (текст) & да & нет \\
\hline
AUHead \cite{lyu2026auhead} & да & да (AU) & да & нет \\
\hline
EmoDiffTalk \cite{liu2026emodifftalk} & да & да (AU) & да & нет \\
\hline
EMOCA \cite{danecek2022emoca} & - & - & да & нет \\
\hline
\textbf{Настоящее исследование} & \textbf{нет} & \textbf{нет} & \textbf{да} & \textbf{да} \\
\hline
\end{tabular}
\end{table}

Анализ сравнительной таблицы выявляет три обстоятельства, существенных для постановки исследования. Во-первых, ни один из рассмотренных подходов не сочетает свойство надстраиваемости к существующим опорным архитектурам с независимой внешней валидацией результата. Во-вторых, подходы, обеспечивающие эмоциональную согласованность без явных меток на этапе инференса, требуют замены базовой архитектуры, что несовместимо с задачей адаптации уже развёрнутых систем. В-третьих, систематический контроль риска цикличности оценки в литературе не реализован: даже работы, явно адресующие эмоциональную согласованность, оценивают результат теми же моделями, что использованы в обучающей сигнальной цепи.

На основании выявленных обстоятельств формулируется \textbf{теоретический пробел}: в литературе отсутствует функция потерь, реализуемая как надстройка над опорными архитектурами Wav2Lip и SadTalker и обеспечивающая повышение эмоциональной согласованности с верификацией результата независимым внешним классификатором, не задействованным в обучающей цепи.

Сужение пробела до надстраиваемого подхода для двух конкретных опорных архитектур мотивировано как теоретически, так и практически. Теоретически такое сужение позволяет сформулировать опровержимую гипотезу с измеримыми порогами, что соответствует требованиям к научной корректности. Практически - уже развёрнутые системы цифровых аватаров используют именно опорные архитектуры рассматриваемых классов, и переход на специализированные эмоциональные модели (EmoTalk, EmotiveTalk и аналогичные) требует существенных архитектурных и инфраструктурных изменений. Снижение барьера адаптации существующих систем к задаче эмоционально-согласованного синтеза составляет практическую значимость настоящего исследования.

Сформулированный теоретический пробел естественным образом ведёт к гипотезам H1 (для опорной архитектуры Wav2Lip) и H2 (для опорной архитектуры SadTalker), приведённым во Введении. Ожидается, что предлагаемая надстраиваемая функция потерь обеспечит статистически значимое улучшение F1-меры эмоциональной согласованности при сохранении LSE-C в пределах допустимого ограничения, причём указанный эффект подтвердится не только по внутренним метрикам, но и при оценке независимым внешним классификатором. Подтверждение или опровержение указанных гипотез составляет предмет эмпирической части настоящей работы.

\subsection{Уточнение позиции работы относительно смежных направлений}

Для полноты теоретической рамки настоящий подраздел уточняет позицию работы относительно нескольких смежных направлений, не охваченных явно в предыдущих обзорах, но методологически связанных с предложенным подходом.

\textbf{Дистилляция знаний между модальностями.} Концептуально победившая в абляции конфигурация (CE+KL) представляет собой частный случай кросс-модальной дистилляции знаний: распределение классов, предсказанное аудиоэнкодером по аудиосигналу, выступает в роли <<мягкой метки>> (soft label) для видеоканала. Классическая работа \cite{hinton2015distillation} формализовала дистилляцию как перенос распределения учителя на ученика через KL-дивергенцию с температурой; CRD \cite{tian2020crd} и SimKD \cite{chen2022simkd} обобщают подход на контрастивные постановки. Применение дистилляции в кросс-модальной постановке ранее изучалось в задачах распознавания --- от аудио к визуальному каналу для задач распознавания речи \cite{afouras2018}, от RGB к глубинному каналу для классификации действий \cite{garcia2018}, --- однако в задаче генерации говорящих лиц дистилляция эмоционального распределения от аудио к видео в режиме надстраиваемого дообучения в литературе системно не исследована. Эмпирический выигрыш CE+KL над чисто супервизионной формулировкой CE-only в проведённой абляции согласуется с известным результатом \cite{hinton2015distillation}: мягкая метка содержит больше информации о структуре пространства классов, чем жёсткая, и при ограниченной обучающей выборке обеспечивает лучшее обобщение.

\textbf{Параметро-эффективный файнтюнинг.} Предложенная функция потерь применяется в режиме обновления всех параметров декодера Wav2Lip и параметров проекций; альтернативой является параметро-эффективный файнтюнинг (LoRA \cite{hu2022lora}, адаптерные модули \cite{houlsby2019adapters}), при котором обновляется только небольшое подмножество параметров. Данное направление выходит за пределы настоящего эксперимента, но методологически совместимо с предложенной функцией потерь и составляет одно из направлений дальнейших работ.

\textbf{Аудио-производные эмбеддинги для управления генерацией.} В классе работ по синтез речи (text-to-speech) и преобразование речи (speech-to-speech) синтезу применяется управление эмоциональной выразительностью через производные аудио-эмбеддинги (представления стиля (style embeddings)) без явной разметки эмоций \cite{wang2018tacotrongst}. Аналогичный принцип в задаче генерации говорящих лиц реализован в \cite{shen2024emotalker}, однако с переобучением архитектуры с нуля. Концептуально CE+KL-сигнал в настоящей работе выполняет ту же функцию --- перенос аудио-производного представления на видеоканал --- но без модификации опорной архитектуры.

\textbf{Регуляризация через выровненность представлений.} Подходы, в которых регуляризатор обеспечивает близость представлений в общем латентном пространстве, представлены в работах по мультимодальной классификации \cite{tsai2019mult} и аудиовизуальному распознаванию \cite{zhang2024}. От указанных работ настоящее исследование отличает применение регуляризатора в задаче генерации (а не классификации) и контроль через категориальную метку эмоции, а не через парные синхронные представления.

\subsection{Объём литературного обзора и принципы отбора источников}

Литературная база настоящего раздела включает не менее 30 первоисточников, отнесённых к следующим тематическим группам: (1) психологические основы классификации эмоций; (2) архитектурные семейства моделей синтеза говорящих лиц (Wav2Lip и его модификации, SadTalker, NeRF-семейство, диффузионные подходы); (3) подходы к интеграции эмоционального компонента в синтез говорящих лиц; (4) предобученные аудио- и видеоэнкодеры и кросс-модальные представления; (5) корпуса с эмоциональной разметкой; (6) контрастивное обучение и дистилляция знаний; (7) метрики оценки качества генерации и проблема цикличности. Принцип отбора --- релевантность к одному из трёх ключевых аспектов работы (опорная архитектура, эмоциональный компонент, протокол оценки) и принадлежность к рецензируемым публикациям профильных конференций (CVPR, ICCV, ECCV, NeurIPS, ICLR, ACL, ICASSP, INTERSPEECH) либо к признанным журналам по машинному обучению. Конкретный отбор не претендует на исчерпывающий охват области, что было бы избыточно для решения сформулированной задачи; вместо этого отобранные источники последовательно реконструируют логику движения от широкого контекста (эмоциональные модели в психологии) к узкой постановке (надстраиваемая функция потерь для двух конкретных опорных архитектур) и формируют обоснование как для гипотез, так и для протокола оценки.


\section{Методология и методика исследования}

\subsection{Общая структура методики}

Методика исследования организована как последовательность пяти взаимосвязанных этапов, выстроенных от подготовки данных и опорных моделей через формализацию функции потерь к экспериментальной проверке гипотез на двух опорных архитектурах. Каждый этап реализован отдельным ноутбуком в открытом репозитории; нумерация ноутбуков (01-11) соответствует порядку выполнения в общем пайплайне:
\begin{enumerate}
\item \textbf{Этап 1 - подготовка данных и разбиение по актёрам:} [01\_data\_preprocessing.ipynb](01\_data\_preprocessing.ipynb).
\item \textbf{Этап 2 - файнтюнинг аудио- и видеоэнкодеров эмоций} (опорные модели функции потерь) и установление <<потолка>> разделимости: [02\_train\_emotion\_encoders.ipynb](02\_train\_emotion\_encoders.ipynb), [03\_emotion\_utils.ipynb](03\_emotion\_utils.ipynb), [04\_encoder\_ceiling.ipynb](04\_encoder\_ceiling.ipynb).
\item \textbf{Этап 3 - скрининг внешних классификаторов и подготовка протокола внешней оценки:} [05\_external\_classifier\_screening.ipynb](05\_external\_classifier\_screening.ipynb).
\item \textbf{Этап 4 - абляционный анализ компонент функции потерь:} [06\_wav2lip\_loss\_ablation.ipynb](06\_wav2lip\_loss\_ablation.ipynb) для Wav2Lip и [09\_sadtalker\_loss\_ablation.ipynb](09\_sadtalker\_loss\_ablation.ipynb) для SadTalker.
\item \textbf{Этап 5 - файнтюнинг опорных архитектур с применением победившей в абляции функции потерь и проверка гипотез:} [07\_wav2lip\_finetune\_H1.ipynb](07\_wav2lip\_finetune\_H1.ipynb), [08\_wav2lip\_external\_evaluation.ipynb](08\_wav2lip\_external\_evaluation.ipynb), [10\_sadtalker\_finetune\_H2.ipynb](10\_sadtalker\_finetune\_H2.ipynb), [11\_bootstrap\_analysis.ipynb](11\_bootstrap\_analysis.ipynb).
\end{enumerate}

Эксперимент в строгом смысле составляет этапы 4 и 5: на этапе 4 непосредственно манипулируется состав функции потерь при фиксированной архитектуре и данных, на этапе 5 - весовой коэффициент функции потерь при фиксированной её формулировке и опорной архитектуре. Все остальные этапы являются подготовительными или валидационными.

\subsection{Подготовка данных и протокол разбиения (этап 1)}

В качестве основного корпуса используется RAVDESS (Ryerson Audio-Visual Database of Emotional Speech and Song) \cite{livingstone2018ravdess}, публично доступный по лицензии CC BY-NA-SC 4.0. Из исходного корпуса извлекаются 1380 аудиовизуальных сэмплов со всеми восемью эмоциональными классами; из них формируется четырёхклассовое подмножество объёмом 736 сэмплов (по 184 на класс): \textit{happy}, \textit{sad}, \textit{angry}, \textit{disgust}. Обоснование выбора четырёх классов приведено в теоретическом разделе. Все 24 актёра RAVDESS разбиты на три непересекающиеся группы по идентификаторам субъектов (\textit{actor-aware split}), что исключает утечку признаков говорящего между сплитами: 19 актёров в обучающей выборке (544 сэмпла), 2-3 актёра в валидационной (96 сэмплов) и 2-3 актёра в тестовой (96 сэмплов, по 24 на каждую из четырёх эмоций).

\textit{Видеообработка.} Каждое исходное видео декодируется на 25~fps; центрированный лицевой кроп извлекается детектором лиц на основе MediaPipe Face Mesh; кадры приводятся к разрешению $96 \times 96$ для совместимости с архитектурой Wav2Lip и сохраняются в виде numpy-тензоров. Для видеоэнкодера эмоций (TimeSformer) дополнительно формируется тензор из 8 кадров, равномерно отобранных по длительности клипа, в разрешении $224 \times 224$ согласно входным требованиям модели. Покадровый успех детекции лица составил 98,9~\% на обучающей выборке и 99,0~\% на валидационной и тестовой; сэмплы с долей успешной детекции ниже 90~\% исключаются из дальнейшей обработки.

\textit{Аудиообработка.} Каждый аудиосигнал ресэмплируется к 16~kHz, моно. Длительность аудиоклипов варьируется от 3,5 до 3,8 секунд (среднее по сплитам). Для совместимости с lipsync-экспертом SyncNet строится мел-спектрограмма с параметрами, унаследованными от исходного пайплайна Wav2Lip: число мел-каналов 80, длина окна 800 семплов, шаг 200 семплов, что соответствует 16 спектральным шагам на каждый видеокадр при 25~fps. Полученная спектрограмма сохраняется как тензор формата $80 \times 16 \times T$, где $T$ --- число видеокадров клипа.

\textit{Метаданные.} Для каждого сэмпла формируется JSON-запись \texttt{metadata.json}, содержащая: идентификатор актёра, индекс эмоции, индекс эмоциональной интенсивности (low/normal), индекс заявления, индекс повтора, путь к аудиофайлу, путь к видеотензору, путь к мел-спектрограмме, длительность клипа, статус детекции лица, принадлежность к сплиту (train/val/test). Все последующие этапы пайплайна загружают данные исключительно через \texttt{metadata.json}, что обеспечивает единообразие и воспроизводимость.

\subsection{Опорные энкодеры эмоций (этап 2)}

Аудиоэнкодер и видеоэнкодер обучаются раздельно - как 4-классовые классификаторы эмоций - с целью получения двух независимых каналов <<эмоционального сигнала>>, выходы которых используются на этапе 5 в качестве целевых распределений в кросс-модальной функции потерь.

\textit{Гиперпараметрический перебор.} Для аудиомодальности рассмотрено 11 конфигураций (варианты \texttt{wav2vec2-base}, \texttt{wav2vec2-large}, \texttt{hubert-large} и SUPERB-вариантов с разными темпами обучения, разными окнами и режимами заморозки/класс-балансировки); для видео - 11 конфигураций \texttt{timesformer-base-finetuned-k400} с числом кадров 4/8/12/16 и темпами обучения от $1\!\cdot\!10^{-5}$ до $5\!\cdot\!10^{-5}$. Все запуски - 30 эпох с ранней остановкой по валидационной F1, оптимизатор AdamW, scheduler CosineAnnealingLR, batch-size 8, label smoothing 0,1, weight decay 0,01, seed 42. Логирование - W\&B (\texttt{uncanny-valley-encoders-4emo}). Полная таблица результатов перебора приведена в разделе результатов.

\textit{Финальный выбор.} По итогам перебора отобраны:
\begin{itemize}
\item \textbf{Аудиоэнкодер:} \texttt{superb/wav2vec2-base-superb-er}, темп обучения $3\!\cdot\!10^{-5}$, длина окна по умолчанию (валидационная макро-F1 0,826).
\item \textbf{Видеоэнкодер:} \texttt{facebook/timesformer-base-finetuned-k400}, 8 кадров на сэмпл, темп обучения $3\!\cdot\!10^{-5}$ (валидационная макро-F1 0,790).
\end{itemize}

Обе модели на этапах 3 и 5 \textit{зафиксированы} (\texttt{requires\_grad=False}); их выходы (4-мерные логиты эмоций) используются исключительно как сигнал в функции потерь, без дальнейшего обновления весов.

\subsection{Формулировка кросс-модальной функции потерь}

Функция потерь, дополняющая стандартную реконструкцию L1 при файнтюнинге опорной архитектуры $\mathcal{G}$, имеет общий вид
\begin{equation}\label{eq:totalloss}
\mathcal{L} = \mathcal{L}_{\text{recon}} + s \cdot \big(\lambda_{\text{ce}}\,\mathcal{L}_{\text{CE}} + \lambda_{\cos}\,\mathcal{L}_{\cos} + \lambda_{\text{KL}}\,\mathcal{L}_{\text{KL}}\big),
\end{equation}
где
\begin{itemize}
\item $\mathcal{L}_{\text{recon}} = \|\hat{V} - V\|_1$ - стандартная L1-реконструкция нижней половины кадра, унаследованная от Wav2Lip;
\item $\mathcal{L}_{\text{CE}} = \mathrm{CE}(\phi_v(\hat{V}),\, y)$ - перекрёстная энтропия между логитами видеоэнкодера на сгенерированном кадре и истинной меткой эмоции $y$;
\item $\mathcal{L}_{\cos} = 1 - \cos\!\bigl(\phi_v(\hat V),\, \phi_a(A)\bigr)$ - косинусное сближение признаков сгенерированного видео и аудиосигнала;
\item $\mathcal{L}_{\text{KL}} = \tau^2\,\mathrm{KL}\!\bigl(\sigma(\phi_v(\hat V)/\tau)\,\|\,\sigma(\phi_a(A)/\tau)\bigr)$ - KL-дистилляция, переносящая мягкое распределение классов с аудиоканала на видеоканал, температура $\tau=2{,}0$;
\item $s$ - общий масштабный коэффициент, $\lambda_{*}$ - покомпонентные веса, выбираемые в абляции.
\end{itemize}

Выбор именно такой формы продиктован тремя соображениями. Во-первых, формулировка \eqref{eq:totalloss} остаётся \textit{надстраиваемой}: она не требует изменений ни в архитектуре $\mathcal{G}$, ни в её прямом проходе. Во-вторых, набор компонент покрывает три качественно различных режима выравнивания: явная супервизия по меткам ($\mathcal{L}_{\text{CE}}$), бесмаркерное выравнивание представлений ($\mathcal{L}_{\cos}$) и непрерывный перенос распределения с одной модальности на другую ($\mathcal{L}_{\text{KL}}$). В-третьих, оптимальное сочетание компонент не известно a priori и устанавливается эмпирически на этапе абляции.

\subsection{Абляционный анализ функции потерь (этап 4, проверка корректности формулировки)}

Цель этапа - эмпирически выбрать состав функции потерь, который при минимальной деградации L1-реконструкции даёт наибольший прирост F1 на валидационной выборке. Поскольку речь о выборе формулы, а не порога, испытываются пять конфигураций:
\begin{enumerate}
\item \textbf{abl-baseline} - $\lambda_{\text{ce}}=\lambda_{\cos}=\lambda_{\text{KL}}=0$ (только реконструкция).
\item \textbf{abl-cos-only} - $\lambda_{\cos}=0{,}2$, остальные веса нулевые. Контроль на коллапс наивной косинусной формулировки в отсутствие негативных пар.
\item \textbf{abl-ce-only} - $\lambda_{\text{ce}}=0{,}05$.
\item \textbf{abl-ce-cos} - $\lambda_{\text{ce}}=0{,}05$, $\lambda_{\cos}=0{,}2$.
\item \textbf{abl-ce-kl} - $\lambda_{\text{ce}}=0{,}05$, $\lambda_{\text{KL}}=0{,}1$, $\tau=2{,}0$.
\end{enumerate}

Все конфигурации обучаются на одном и том же seed (42) на 70 эпохах с ранней остановкой по 8 эпохам без улучшения валидационной F1 (для эмоциональных конфигураций) либо валидационной L1 (для baseline). Опорная архитектура - Wav2Lip; все слои трансформера и SyncNet-эксперт заморожены, обновляются только параметры аудио- и видеоэнкодеров и декодера Wav2Lip. Полная техническая спецификация запусков и логи доступны в проекте \texttt{uncanny-valley-wav2lip-ablation} в W\&B.

Для каждой конфигурации регистрируются: динамика валидационной L1 и F1 по эпохам; финальное значение L1, F1, accuracy на валидационной выборке; статистическая значимость отличия от baseline (для F1 - критерий Макнемара; для L1 - парный $t$-критерий и критерий Уилкоксона). Победившая по F1 конфигурация переходит на этап 5 для проверки гипотез на тестовой выборке.

\subsection{Поиск весового коэффициента и проверка H1 (этап 5, Wav2Lip)}

Победившая в абляции комбинация (CE+KL) применяется к Wav2Lip с переменным масштабом $s\in\{0{,}1;\,0{,}25;\,0{,}4;\,0{,}5\}$. Для каждой точки сохраняется чекпойнт по максимуму валидационной F1. Правило отбора финальной модели для проверки H1 - максимум F1 среди конфигураций, не выходящих за полосу <<+150~\% L1-реконструкции от baseline>>. Указанное послабление полосы $L_1$ обусловлено тем, что критерий H1 формулируется не на $L_1$, а на LSE-C тестовой выборки; полоса $L_1$ выступает санитарным фильтром против заведомо испорченных конфигураций, а финальный вердикт по сохранности артикуляции выносится по LSE-C на отложенной тестовой выборке.

После отбора финальной модели проводится сравнение с базовой моделью (\texttt{wav2lip-baseline}, нулевой эмоциональный лосс) на тестовой выборке RAVDESS ($n=96$, по 24 сэмпла на эмоцию):
\begin{itemize}
\item категориальная метрика эмоциональной согласованности по внутреннему видеоэнкодеру (accuracy, макро-F1, per-emotion P/R/F1, матрица ошибок);
\item LSE-C / LSE-D, вычисляемые предобученным lipsync-экспертом SyncNet;
\item L1-реконструкция как контрольная величина.
\end{itemize}
Статистическая значимость: критерий Макнемара для F1 (с подсчётом дискордантных пар $n_{01}$, $n_{10}$, $\chi^2$, $p$); парный $t$-критерий и критерий Уилкоксона для L1 и LSE-C. Решение по H1 принимается при одновременном выполнении двух критериев: $\Delta\text{F1} \geq +10$~\% при $p < 0{,}05$ (Макнемар) и $|\Delta\text{LSE-C}| \leq 2$~\% либо $p \geq 0{,}05$ (парный $t$-критерий).

\subsection{Внешняя валидация и контроль цикличности оценки (этапы 3, 5)}

Поскольку и функция потерь, и первичная метрика на этапе 5 опираются на один и тот же видеоэнкодер $\phi_v$, оценка только по $\phi_v$ создаёт риск кругового рассуждения. Для устранения риска проведён скрининг внешних классификаторов выражений лица в [05\_external\_classifier\_screening.ipynb](05\_external\_classifier\_screening.ipynb): из шести кандидатов один (\texttt{prithivMLmods/Facial-Emotion-Detection-SigLIP2}) исключён из-за отсутствия класса \textit{disgust} в его выходном пространстве; пять оставшихся прогнаны на тестовой выборке RAVDESS в режиме классификации центральных кадров реальных видео, и для основной внешней оценки отобраны два классификатора с максимальной макро-F1 на четырёх целевых эмоциях - \texttt{motheecreator/vit-Facial-Expression-Recognition} и \texttt{Rajaram1996/FacialEmoRecog}; \texttt{trpakov/vit-face-expression} удерживается как контрольный классификатор для отдельных подвыборок.

На этапе 5 финальные чекпойнты Wav2Lip применяются к тестовому подмножеству, генерируется выход на каждый сэмпл, центральный кадр из нижней половины каждого сгенерированного видео классифицируется каждым внешним классификатором, и подсчитывается макро-F1, accuracy, per-emotion F1 и матрица ошибок. Для каждой пары <<конфигурация, внешний классификатор>> рассчитывается $\Delta$F1 относительно базовой модели и проводится критерий Макнемара. В качестве диагностической метрики фиксируется \textit{межклассификаторное согласие} - доля конфигураций, для которых направление $\Delta$F1 совпадает между двумя различающимися по архитектуре и обучающему корпусу классификаторами.

В качестве вспомогательного аудио-канала проверки фиксируется регрессия arousal/dominance/valence по модели \texttt{audeering/wav2vec2-large-robust-12-ft-emotion-msp-dim} (MSP-Podcast), что позволяет соотнести предсказанные эмоциональные категории с непрерывными размерными признаками речи и тем самым отделить случайные совпадения меток от содержательно согласованного эмоционального тона.

\subsection{Перенос на SadTalker и проверка H2 (этап 5)}

Опорная архитектура SadTalker \cite{zhang2023sadtalker} отличается от Wav2Lip тем, что генерирует параметры 3DMM-геометрии лица, а не пиксели нижней половины кадра. Для применения той же функции потерь \eqref{eq:totalloss} в её декомпозиции <<ExpNet $\to$ FaceRender>> промежуточный сигнал отбирается на этапе синтеза кадров: после рендеринга применяется тот же видеоэнкодер $\phi_v$ к сгенерированному кадру, и далее используются $\mathcal{L}_{\text{CE}}$ и $\mathcal{L}_{\text{KL}}$. Файнтюнинг проводится на тех же 544 обучающих сэмплах RAVDESS, что и Wav2Lip; сохраняется тот же опорный аудиоэнкодер. Для скрининга масштаба $s$ в [10\_sadtalker\_finetune\_H2.ipynb](10\_sadtalker\_finetune\_H2.ipynb) проверены конфигурации $s=0{,}25$ и $s=1{,}0$; в [09\_sadtalker\_loss\_ablation.ipynb](09\_sadtalker\_loss\_ablation.ipynb) дополнительно проверена конфигурация \textit{abl-ce-cos}, заявленная как лучшая по средней внешней $\Delta$F1 в эксперименте Wav2Lip (см. результаты).

Внешняя оценка SadTalker отличается от Wav2Lip следующим: вместо классификации центральных кадров библиотечных тестовых видео внешний классификатор применяется к рендеренным видео (один кадр из каждых четырёх) с агрегированием по голосованию большинства; объём финальной выборки - 24 видео на конфигурацию (по 6 на эмоцию), что обусловлено вычислительной стоимостью рендеринга. Указанное ограничение явно отражено в разделе результатов и в обсуждении ограничений; малый размер выборки исключает чувствительное обнаружение слабых эффектов.

Решение по H2: $\Delta\text{F1} \geq +8$~\% по обоим внешним классификаторам с положительным согласованием направления и $p < 0{,}05$ по критерию Макнемара хотя бы для одного из них.

\subsection{Технические детали обучения (общие для этапов 2, 4, 5)}

Все запуски в работе используют единые технические настройки, перечисленные ниже; локальные отклонения от настроек явно указаны в соответствующих подразделах.

\textit{Аппаратная база.} Обучение проводится на Google Colab с одним графическим ускорителем NVIDIA Tesla T4 (16~GB VRAM) либо A100 (40~GB) в зависимости от доступности. Локальные эксперименты по внешней оценке проведены на NVIDIA GeForce RTX 3090 (24~GB VRAM). CUDA версия 12.4, драйвер 580.142.

\textit{Программная база.} Python 3.10, PyTorch 2.1, Hugging Face Transformers 4.40, scikit-learn 1.3, scipy 1.11, librosa 0.10, torchaudio 2.1. Для логирования и визуализации --- Weights \& Biases (\texttt{wandb} 0.16). Полный закреплённый список зависимостей входит в первую ячейку каждого ноутбука.

\textit{Воспроизводимость.} Фиксированный seed=42 устанавливается для модулей \texttt{random}, \texttt{numpy.random}, \texttt{torch.manual\_seed}, \texttt{torch.cuda.manual\_seed\_all}; режим \texttt{torch.backends.cudnn.deterministic = True} активирован для всех запусков. DataLoader использует \texttt{num\_workers=2} с фиксированным \texttt{worker\_init\_fn} для повторяемости перемешивания.

\textit{Оптимизатор и регуляризация.} AdamW с темпом обучения, варьируемым по конфигурации, и weight\_decay~0,01. Scheduler --- CosineAnnealingLR с warmup в 5 эпох (только для финальных файнтюнингов на этапе 5; для этапов 2--3 warmup не применялся). Label smoothing~0,1 для CE-членов в этапах 2 и 3. Mixed precision (\texttt{torch.amp}) применяется для аудиоэнкодеров на крупных архитектурах (HuBERT-Large, Wav2Vec2-Large) для снижения требований к памяти.

\textit{Ранняя остановка.} Для всех запусков применяется ранняя остановка по критерию <<нет улучшения по выбранной метрике в течение $K$ эпох>>; $K=8$ для этапов 2 и 4, $K=8$ для этапа 5. Метрика отбора чекпойнта: для этапов 2 и 5 --- максимум валидационной макро-F1; для этапа 4 в baseline-конфигурации (без эмоционального члена) --- минимум валидационной L1, в эмоциональных конфигурациях --- максимум валидационной макро-F1.

\subsection{Алгоритмическая схема функции потерь (псевдокод)}

Для устранения двусмысленностей в реализации функции потерь приведём её алгоритмическую схему. На каждом обучающем шаге для батча из $B$ сэмплов вычисляются:

\begin{enumerate}
\item Прогон опорной архитектуры (Wav2Lip): $\hat{V} \leftarrow \mathcal{G}(A_{\text{mel}}, V_{\text{ref}})$, где $A_{\text{mel}}$ --- мел-спектрограмма аудиоклипа, $V_{\text{ref}}$ --- референсные кадры с открытым лицом, $\hat{V}$ --- сгенерированная нижняя половина лица.
\item Реконструкционный член: $\mathcal{L}_{\text{recon}} \leftarrow \|\hat{V} - V_{\text{gt}}\|_1$ по нижней половине кадра.
\item Прогон зафиксированного видеоэнкодера на сгенерированных кадрах: $\mathbf{l}_v \leftarrow \phi_v(\hat{V}) \in \mathbb{R}^{B \times 4}$ (4-мерные логиты эмоций).
\item Прогон зафиксированного аудиоэнкодера на исходном аудио: $\mathbf{l}_a \leftarrow \phi_a(A) \in \mathbb{R}^{B \times 4}$, $\mathbf{l}_a \leftarrow \mathbf{l}_a.\mathrm{detach}()$ (отсечение градиента).
\item Перекрёстная энтропия по меткам: $\mathcal{L}_{\text{CE}} \leftarrow \mathrm{CE}(\mathbf{l}_v, y)$.
\item KL-дистилляция при температуре $T$: $\mathcal{L}_{\text{KL}} \leftarrow T^2 \cdot \mathrm{KL}\bigl(\mathrm{logsoftmax}(\mathbf{l}_v / T)\,\|\,\mathrm{softmax}(\mathbf{l}_a / T)\bigr)$.
\item Косинусный член (если активен): $\mathbf{p}_v \leftarrow W_v \mathbf{l}_v$, $\mathbf{p}_a \leftarrow W_a \mathbf{l}_a$ (проекции в общее пространство), $\mathcal{L}_{\cos} \leftarrow 1 - \cos(\mathbf{p}_v, \mathbf{p}_a)$, усреднено по батчу.
\item Итоговая функция потерь: $\mathcal{L} \leftarrow \mathcal{L}_{\text{recon}} + s \cdot (\lambda_{\text{ce}}\,\mathcal{L}_{\text{CE}} + \lambda_{\text{KL}}\,\mathcal{L}_{\text{KL}} + \lambda_{\cos}\,\mathcal{L}_{\cos})$.
\item Обратное распространение градиента: только параметры $\mathcal{G}$ и проекций $W_v, W_a$ получают градиент; $\phi_v$ и $\phi_a$ заморожены.
\item Шаг оптимизатора с ограничением нормы градиента \texttt{clip\_grad\_norm}~$=2{,}0$.
\end{enumerate}

Полная реализация приведена в Приложении~Б; W\&B-логи каждого запуска фиксируют динамику всех компонент функции потерь по эпохам и шагам, что обеспечивает аудит и воспроизводимость.

\subsection{Воспроизводимость и доступность артефактов}

Все ноутбуки публично доступны в репозитории (\url{https://github.com/Katrin-Pochtar/The-Uncanny-Valley}); фиксированный seed=42 для всех численных операций (Python \texttt{random}, NumPy, PyTorch, CUDA); версии библиотек закреплены в \texttt{requirements} и в выводе \texttt{pip list} первой ячейки каждого ноутбука; гиперпараметры и метрики записаны в проекты W\&B \texttt{uncanny-valley-encoders-4emo} (этап~2), \texttt{uncanny-valley-wav2lip-ablation} (этап~3), \texttt{uncanny-valley-wav2lip-4emo} (этап~5, Wav2Lip), \texttt{uncanny-valley-sadtalker-ablation-4emo} (этап~5, SadTalker). Чекпойнты опорных энкодеров, абляции и финальных моделей сохранены в локальных каталогах, структура которых описана в README репозитория. Реестр артефактов с явной связью <<артефакт $\to$ шаг методики $\to$ результат>> приведён в Приложении~А.

\section{Результаты}

\subsection{Подбор опорных энкодеров эмоций}

В рамках этапа~2 проведено 22 запуска (11 аудиоконфигураций и 11 видеоконфигураций); полная таблица результатов перебора приведена в Приложении~В, лучшие конфигурации обоих модальностей представлены в табл.~\ref{tab:enc_sweep}.

\textit{Динамика обучения и кривые валидации} для всех конфигураций зафиксированы в проекте W\&B \texttt{uncanny-valley-encoders-4emo}. Типичная динамика для победившей аудиоконфигурации (\texttt{wav2vec2-base-superb-er}, lr~$3\!\cdot\!10^{-5}$): валидационная макро-F1 достигает 0,80 к 7 эпохе, далее медленный рост до 0,826 на 17 эпохе с активацией ранней остановки на 23 эпохе. Для видеоконфигурации (\texttt{timesformer-base-finetuned-k400}, 8 кадров, lr~$3\!\cdot\!10^{-5}$): валидационная макро-F1 достигает 0,76 к 7 эпохе и стабилизируется в диапазоне 0,77--0,79 с активацией ранней остановки на 24--26 эпохе.

\begin{center}
\textit{[\textbf{W\&B figure}: training-curves-encoders --- две панели (audio/video) с кривыми валидационной макро-F1 vs эпохи для топ-5 конфигураций каждой модальности; источник: проект \texttt{uncanny-valley-encoders-4emo}, виджет <<Compare runs>>]}
\end{center}

\begin{table}[H]
\caption{Лучшие конфигурации энкодеров эмоций (валидационная макро-F1, RAVDESS, 4 эмоции)}
\label{tab:enc_sweep}
\centering
\small
\begin{tabular}{|p{5.5cm}|c|c|c|}
\hline
\textbf{Конфигурация} & \textbf{Модальность} & \textbf{Лучшая Val~F1} & \textbf{Финал} \\
\hline
\texttt{wav2vec2-base-superb-er}, lr $3\!\cdot\!10^{-5}$ & аудио & 0,826 & да \\
\hline
\texttt{hubert-large-ll60k}, lr $2\!\cdot\!10^{-5}$ & аудио & 0,804 & - \\
\hline
\texttt{wav2vec2-large}, lr $2\!\cdot\!10^{-5}$ & аудио & 0,801 & - \\
\hline
\texttt{timesformer-base}, 8 кадров, lr $3\!\cdot\!10^{-5}$ & видео & 0,790 & да \\
\hline
\texttt{timesformer-base}, 16 кадров, no-freeze, lr $3\!\cdot\!10^{-5}$ & видео & 0,781 & - \\
\hline
\texttt{timesformer-base}, 16 кадров, lr $5\!\cdot\!10^{-5}$ & видео & 0,780 & - \\
\hline
\end{tabular}
\end{table}

На независимой тестовой выборке отобранные энкодеры демонстрируют макро-F1 0,739 (аудио) и 0,876 (видео) ([04\_encoder\_ceiling.ipynb](04\_encoder\_ceiling.ipynb)). Поэмоциональная разбивка приведена в табл.~\ref{tab:enc_ceiling}; видеоэнкодер уверенно классифицирует \textit{happy}, \textit{angry} и \textit{disgust} (F1 $\geq$ 0,90) и хуже работает на \textit{sad} (F1 0,72), что согласуется с известной из литературы низкой различимостью <<печали>> по статичной мимике.

\begin{table}[H]
\caption{<<Потолок>> разделимости энкодеров: per-emotion F1 на тестовой выборке RAVDESS ($n=96$)}
\label{tab:enc_ceiling}
\centering
\small
\begin{tabular}{|p{4cm}|c|c|c|c|c|}
\hline
\textbf{Энкодер} & \textbf{Macro F1} & \textbf{F1 happy} & \textbf{F1 sad} & \textbf{F1 angry} & \textbf{F1 disgust} \\
\hline
Аудио (\texttt{wav2vec2-base-superb-er}) & 0,739 & 0,667 & 0,760 & 0,721 & 0,810 \\
\hline
Видео (\texttt{timesformer-base}, 8 fr.) & 0,876 & 0,906 & 0,718 & 0,923 & 0,958 \\
\hline
\end{tabular}
\end{table}

\subsection{Скрининг внешних классификаторов}

В рамках этапа~3 шесть кандидатов рассмотрены и пять прогнаны на тестовой выборке RAVDESS ([05\_external\_classifier\_screening.ipynb](05\_external\_classifier\_screening.ipynb)); один (\texttt{prithivMLmods/Facial-Emotion-Detection-SigLIP2}) исключён из-за отсутствия класса \textit{disgust}. Результаты приведены в табл.~\ref{tab:ext_screen}.

\begin{table}[H]
\caption{Скрининг внешних классификаторов выражений лица на тестовой выборке RAVDESS, $n=96$}
\label{tab:ext_screen}
\centering
\small
\begin{tabular}{|p{6cm}|c|c|c|c|c|}
\hline
\textbf{Классификатор} & \textbf{Macro F1} & \textbf{F1 happy} & \textbf{F1 sad} & \textbf{F1 angry} & \textbf{F1 disgust} \\
\hline
\texttt{motheecreator/vit-FER} & \textbf{0,476} & 0,738 & 0,133 & 0,410 & 0,621 \\
\hline
\texttt{trpakov/vit-face-expression} & 0,464 & 0,814 & 0,364 & 0,677 & 0,000 \\
\hline
\texttt{Rajaram1996/FacialEmoRecog} & \textbf{0,462} & 0,744 & 0,250 & 0,263 & 0,592 \\
\hline
\texttt{RickyIG/emotion\_face} & 0,432 & 0,718 & 0,486 & 0,000 & 0,522 \\
\hline
\texttt{dima806/facial\_emotions} & 0,280 & 0,608 & 0,065 & 0,448 & 0,000 \\
\hline
\end{tabular}
\end{table}

Промежуточная макро-F1 внешних классификаторов на чистом RAVDESS (0,28-0,48) существенно ниже <<потолка>> внутреннего видеоэнкодера (0,876), что фиксирует эффект доменного сдвига между обучающими корпусами внешних классификаторов (преимущественно AffectNet/FER+) и студийной актёрской записью RAVDESS. Указанный сдвиг определяет верхнюю границу любого эффекта, регистрируемого внешними классификаторами на сгенерированном видео: даже идеальный генератор не может превзойти внешнюю F1 на реальных кадрах. Для финальной внешней оценки отобраны \texttt{motheecreator/vit-FER} и \texttt{Rajaram1996/FacialEmoRecog}.

\subsection{Абляция функции потерь (этап 4)}

Результаты пятиконфигурационной абляции на валидационной выборке Wav2Lip ($n=96$) представлены в табл.~\ref{tab:ablation}. Победившая по F1 конфигурация - \textbf{abl-ce-kl} (комбинация перекрёстной энтропии и KL-дистилляции).

\begin{table}[H]
\caption{Абляция функции потерь на Wav2Lip, валидационная выборка RAVDESS}
\label{tab:ablation}
\centering
\small
\begin{tabular}{|p{2.5cm}|c|c|c|c|c|c|}
\hline
\textbf{Конфигурация} & $w_{\text{ce}}$ & $w_{\cos}$ & $w_{\text{KL}}$ & \textbf{Best L1} & \textbf{Best F1} & $\Delta$F1 \\
\hline
abl-baseline & 0,00 & 0,0 & 0,00 & 0,0388 & 0,619 & 0,000 \\
\hline
abl-cos-only & 0,00 & 0,2 & 0,00 & 0,0420 & 0,574 & $-0,045$ \\
\hline
abl-ce-only & 0,05 & 0,0 & 0,00 & 0,0619 & 0,696 & $+0,077$ \\
\hline
abl-ce-cos & 0,05 & 0,2 & 0,00 & 0,0654 & 0,655 & $+0,036$ \\
\hline
\textbf{abl-ce-kl} & \textbf{0,05} & \textbf{0,0} & \textbf{0,10} & \textbf{0,0831} & \textbf{0,737} & \textbf{$+0,118$} \\
\hline
\end{tabular}
\end{table}

\textit{Содержательный вывод абляции:} наивная косинусная формулировка без явных негативных пар (abl-cos-only) приводит к \textit{отрицательной} $\Delta$F1 относительно baseline, что эмпирически воспроизводит описанный в литературе риск тривиального коллапса проекций \cite{chen2021simsiam}. Включение явной супервизии в виде перекрёстной энтропии (abl-ce-only) даёт умеренный прирост; добавление к ней KL-дистилляции от аудиоэнкодера к видеоэнкодеру (abl-ce-kl) увеличивает прирост ещё в полтора раза, при этом критерий Макнемара на $n=96$ показывает $p$-значения, не достигающие порога $p<0{,}05$ (для лучшей конфигурации abl-ce-kl: $n_{10}=22$, $n_{01}=14$, $p = 0{,}243$), что обусловлено малым объёмом валидационной выборки. Деградация L1-реконструкции значима ($p < 10^{-9}$ во всех эмоциональных конфигурациях): эмоциональный сигнал смещает декодер в сторону более эмоционально выраженной мимики ценой более грубой попиксельной реконструкции, что и составляет содержание исследуемого компромисса. Финальный вердикт по сохранности артикуляции выносится на тестовой выборке по LSE-C, не по L1.

\begin{center}
\textit{[\textbf{W\&B figure}: ablation-loss-curves --- три панели: (а) валидационная L1-реконструкция, (б) валидационная макро-F1, (в) полная функция потерь $\mathcal{L}_{\text{total}}$ для пяти конфигураций абляции; источник: проект \texttt{uncanny-valley-wav2lip-ablation}, виджет <<Compare runs>>]}
\end{center}

Внешняя оценка той же абляции на тестовой выборке (этап 5, [08\_wav2lip\_external\_evaluation.ipynb](08\_wav2lip\_external\_evaluation.ipynb)) воспроизводит ту же иерархию: abl-cos-only имеет наименьший $\Delta$F1 среди эмоциональных конфигураций, abl-ce-cos и abl-ce-kl --- максимальный; направление эффекта согласовано между двумя независимыми внешними классификаторами для всех восьми эмоциональных конфигураций (межклассификаторное согласие = True). Сводная таблица внешней оценки приведена в табл.~\ref{tab:abl_external}.

\begin{table}[H]
\caption{Внешняя оценка абляции и финальных конфигураций на тестовой выборке Wav2Lip}
\label{tab:abl_external}
\centering
\small
\begin{tabular}{|p{2.5cm}|c|c|c|c|c|}
\hline
\textbf{Конфигурация} & \textbf{Internal val F1} & $\Delta_{\text{mothee.}}$ & $\Delta_{\text{Rajaram.}}$ & \textbf{both\_pos.} & \textbf{Mean $\Delta$} \\
\hline
abl-ce-cos & 0,655 & $+0{,}081$ & $+0{,}051$ & True & $+0{,}066$ \\
\hline
cekl-05 & 0,736 & $+0{,}033$ & $+0{,}089$ & True & $+0{,}061$ \\
\hline
cekl-04 & 0,739 & $+0{,}059$ & $+0{,}039$ & True & $+0{,}049$ \\
\hline
abl-ce-only & 0,696 & $+0{,}049$ & $+0{,}006$ & True & $+0{,}027$ \\
\hline
abl-ce-kl & 0,737 & $+0{,}034$ & $+0{,}013$ & True & $+0{,}024$ \\
\hline
cekl-01 & 0,742 & $+0{,}040$ & $+0{,}007$ & True & $+0{,}023$ \\
\hline
cekl-025 & 0,691 & $+0{,}012$ & $+0{,}009$ & True & $+0{,}010$ \\
\hline
abl-cos-only & 0,574 & $+0{,}006$ & $+0{,}007$ & True & $+0{,}006$ \\
\hline
abl-baseline & 0,619 & $-0{,}001$ & $+0{,}010$ & False & $+0{,}005$ \\
\hline
baseline & 0,624 & 0 & 0 & --- & 0 \\
\hline
\end{tabular}
\end{table}

\textit{Контроль через межклассификаторное согласие.} Все 8 эмоциональных конфигов показывают одновременный положительный $\Delta\text{F1}$ на двух независимых внешних классификаторах; контрольная конфигурация \textit{abl-baseline} (recon-only без эмоционального члена, по дизайну эквивалентная baseline) показывает несущественные смещения и не проходит критерий межклассификаторное согласие, что подтверждает корректность процедуры. Корреляция Спирмена между внутренней валидационной F1 и средним внешним $\Delta\text{F1}$ на тесте составляет $\rho = +0{,}612$ ($p = 0{,}06$ по motheecreator), что свидетельствует о \textit{положительной} переносимости рейтинга конфигураций между внутренним и внешним протоколами; следовательно, выбор конфигурации по внутренней метрике в умеренной степени переносится на внешнюю оценку, что снижает риск кругового рассуждения.

\textit{Контрольное наблюдение по эффекту скрининга внешних.} До скрининга внешних классификаторов в качестве эвалюаторов рассматривались \texttt{dima806/facial\_emotions\_image\_detection} и \texttt{trpakov/vit-face-expression}. На указанной паре ни одна из CE-содержащих конфигураций не давала одновременного положительного $\Delta\text{F1}$ на обоих эвалюаторах: для \texttt{dima806} большинство $\Delta$ положительны, тогда как для \texttt{trpakov} большинство $\Delta$ отрицательны (трактовка: \texttt{trpakov} имеет сильное искажение в сторону \textit{happy}, см. F1\_disgust = 0 и F1\_happy = 0,81 в табл.~\ref{tab:ext_screen}; CE-сигнал перераспределяет предсказания и формально снижает F1 на этом классификаторе из-за роста ошибок типа <<true happy предсказан как другой класс>>). Замена эвалюаторов на отобранные по итогам скрининга \texttt{motheecreator} и \texttt{Rajaram1996} приводит к согласованным положительным $\Delta$ на 8 из 8 эмоциональных конфигураций. Указанное наблюдение составляет самостоятельный методологический результат: вердикт о работоспособности метода эмоциональной супервизии может зависеть от выбора эвалюатора, и предварительный скрининг кандидатов-эвалюаторов на чистом таргет-домене является необходимой частью протокола внешней оценки.

\subsection{H1: проверка на тестовой выборке Wav2Lip}

Победившая в абляции комбинация (CE+KL, $\tau=2{,}0$) применена к Wav2Lip с переменным масштабом $s$ ([07\_wav2lip\_finetune\_H1.ipynb](07\_wav2lip\_finetune\_H1.ipynb)). Сводная таблица валидационных результатов:

\begin{table}[H]
\caption{Поиск масштаба $s$ для CE+KL-функции потерь, валидационная выборка}
\label{tab:scale_sweep}
\centering
\small
\begin{tabular}{|p{2.7cm}|c|c|c|c|c|}
\hline
\textbf{Конфигурация} & $s$ & $w_{\text{ce}}$ & $w_{\text{KL}}$ & \textbf{Best L1} & \textbf{Best F1} \\
\hline
wav2lip-baseline & 0,00 & 0,0000 & 0,000 & 0,0382 & 0,624 \\
\hline
\textbf{wav2lip-cekl-01} & \textbf{0,10} & \textbf{0,0050} & \textbf{0,010} & \textbf{0,0483} & \textbf{0,742} \\
\hline
wav2lip-cekl-04 & 0,40 & 0,0200 & 0,040 & 0,0615 & 0,739 \\
\hline
wav2lip-cekl-05 & 0,50 & 0,0250 & 0,050 & 0,0925 & 0,736 \\
\hline
wav2lip-cekl-025 & 0,25 & 0,0125 & 0,025 & 0,0894 & 0,691 \\
\hline
\end{tabular}
\end{table}

Финальный чекпойнт - \textbf{wav2lip-cekl-01} (максимум валидационной F1 при наименьшей деградации L1; находится в полосе допустимой реконструкции). Сравнение с базовой моделью на тестовой выборке ($n=96$):

\begin{table}[H]
\caption{H1: \textit{wav2lip-baseline} vs \textit{wav2lip-cekl-01}, тестовая выборка RAVDESS}
\label{tab:h1_test}
\centering
\small
\begin{tabular}{|p{4cm}|c|c|c|c|}
\hline
\textbf{Метрика} & \textbf{Baseline} & \textbf{CE+KL} & $\Delta$ & $p$-value \\
\hline
Accuracy (внутр. видеоэнкодер) & 0,635 & 0,719 & $+0,083$ & 0,170 (Mc) \\
\hline
Macro F1 (внутр. видеоэнкодер) & 0,611 & 0,718 & $+0,107$ & 0,170 (Mc) \\
\hline
LSE-C $\uparrow$ & 0,242 & 0,227 & $-0,015$ ($-6{,}1$~\%) & 0,611 ($t$) \\
\hline
LSE-D $\downarrow$ & 0,758 & 0,773 & $+0,015$ & 0,611 ($t$) \\
\hline
L1-реконструкция $\downarrow$ & 0,047 & 0,072 & $+0,025$ & $4{,}4\!\cdot\!10^{-13}$ ($W$) \\
\hline
\end{tabular}
\end{table}

Поэмоциональная разбивка F1 (табл.~\ref{tab:h1_per_emo}) показывает, что наибольший прирост приходится на классы, наиболее <<пересекаемые>> с базовой моделью - \textit{sad} ($+0{,}205$) и \textit{angry} ($+0{,}120$); эффект на \textit{happy} мал ($+0{,}031$) ввиду насыщения базовой модели на этом классе (recall = 1,000).

\begin{table}[H]
\caption{H1: per-emotion F1 на тестовой выборке (Wav2Lip)}
\label{tab:h1_per_emo}
\centering
\small
\begin{tabular}{|c|c|c|c|}
\hline
\textbf{Эмоция} & \textbf{Baseline F1} & \textbf{CE+KL F1} & $\Delta$F1 \\
\hline
happy & 0,762 & 0,792 & $+0,031$ \\
\hline
sad & 0,378 & 0,583 & $+0,205$ \\
\hline
angry & 0,653 & 0,773 & $+0,120$ \\
\hline
disgust & 0,651 & 0,723 & $+0,072$ \\
\hline
\end{tabular}
\end{table}

\textit{Вердикт по H1.} Магнитуда эффекта по F1 удовлетворяет порогу гипотезы ($\Delta\text{F1} = +10{,}7$~\% $\geq +10$~\%); статистическая значимость по критерию Макнемара на $n=96$ не достигнута ($p = 0{,}170$). Изменение LSE-C по абсолютной величине превосходит порог в 2~\% ($\Delta = -6{,}1$~\%), однако не является статистически значимым ($p = 0{,}611$). Согласно сформулированному в методологии правилу решения, $|\Delta\text{LSE-C}|>2$~\% при $p \geq 0{,}05$ удовлетворяет ослабленному критерию (изменение неотличимо от нуля); таким образом, ограничение на сохранность артикуляции \textbf{не нарушено}, тогда как улучшение F1 \textbf{не получило статистического подтверждения} на текущем размере выборки. \textbf{Гипотеза H1 в её строгой формулировке не отвергнута и не подтверждена}; направление эффекта систематично (одинаково на двух внешних классификаторах, см. ниже), но статистическая мощность $n=96$ недостаточна для обнаружения эффекта искомой величины.

\subsection{H1: внешняя валидация}

Внешняя оценка той же пары моделей на двух независимых классификаторах (табл.~\ref{tab:h1_external}) показывает положительный $\Delta$F1 на обоих и направленную согласованность между внешними классификаторами и внутренней метрикой.

\begin{table}[H]
\caption{H1: внешняя валидация на тестовой выборке (Wav2Lip)}
\label{tab:h1_external}
\centering
\small
\begin{tabular}{|p{3.5cm}|c|c|c|c|}
\hline
\textbf{Внешний классификатор} & \textbf{Baseline F1} & \textbf{CE+KL F1} & $\Delta$F1 & McNemar $p$ \\
\hline
\texttt{motheecreator/vit-FER} & 0,303 & 0,343 & $+0,040$ & 0,286 \\
\hline
\texttt{Rajaram1996/FacialEmoRecog} & 0,338 & 0,344 & $+0,007$ & 1,000 \\
\hline
\end{tabular}
\end{table}

Межклассификаторное согласие по направлению эффекта - True для финального чекпойнта \textit{cekl-01}. На альтернативных конфигурациях из той же CE+KL-сем\'ьи достигаются б\'ольшие положительные $\Delta$F1 на обоих внешних классификаторах (\textit{cekl-04}: $+0{,}059$ и $+0{,}039$; \textit{cekl-05}: $+0{,}033$ и $+0{,}089$ соответственно), но при б\'ольшей деградации $L_1$-реконструкции, что и обусловило выбор \textit{cekl-01} как финальной модели по правилу селекции <<максимум F1 при минимуме L1>>. Статистическая значимость на $n=96$ не достигнута ни для одной пары; направление эффекта при этом устойчиво между двумя внешними классификаторами, что снижает риск артефакта одного источника оценки.

\subsection{H1: бутстрап-анализ доверительных интервалов $\Delta\mathrm{F1}$}

Для повышения надёжности вывода критерий Макнемара дополнен парным бутстрап-анализом (10\,000 реплик с возвращением, паросочетание по \texttt{sample\_id}); 95~\%-доверительный интервал на $\Delta\mathrm{F1}$ строится перцентильным методом. Ключевой критерий значимости в бутстрап-парадигме - невключение нуля в доверительный интервал. Полный набор результатов приведён в [11\_bootstrap\_analysis.ipynb](11\_bootstrap\_analysis.ipynb); сводка для отобранных эвалюаторов на тестовой выборке представлена в табл.~\ref{tab:bootstrap_test}.

\begin{table}[H]
\caption{Бутстрап-95~\% CI на $\Delta\mathrm{F1}$ против baseline, тестовая выборка, $n=96$, 10\,000 реплик}
\label{tab:bootstrap_test}
\centering
\small
\begin{tabular}{|p{2.6cm}|p{2.7cm}|c|c|c|}
\hline
\textbf{Внешний} & \textbf{Конфигурация} & $\Delta\mathrm{F1}$ & \textbf{95~\% CI} & \textbf{CI $\not\ni$ 0} \\
\hline
\multirow{4}{*}{Rajaram1996} & cekl-05 & $+0{,}089$ & $[+0{,}004,\,+0{,}176]$ & \textbf{да} \\
\cline{2-5}
& abl-ce-cos & $+0{,}051$ & $[-0{,}040,\,+0{,}134]$ & нет \\
\cline{2-5}
& cekl-04 & $+0{,}039$ & $[-0{,}048,\,+0{,}120]$ & нет \\
\cline{2-5}
& cekl-01 & $+0{,}007$ & $[-0{,}078,\,+0{,}093]$ & нет \\
\hline
\multirow{4}{*}{motheecreator} & abl-ce-cos & $+0{,}082$ & $[-0{,}025,\,+0{,}186]$ & нет \\
\cline{2-5}
& cekl-04 & $+0{,}060$ & $[-0{,}032,\,+0{,}155]$ & нет \\
\cline{2-5}
& cekl-01 & $+0{,}041$ & $[-0{,}052,\,+0{,}138]$ & нет \\
\cline{2-5}
& abl-ce-kl & $+0{,}035$ & $[-0{,}052,\,+0{,}119]$ & нет \\
\hline
\multirow{2}{*}{dima806} & abl-ce-cos & $+0{,}027$ & $[-0{,}032,\,+0{,}083]$ & нет \\
\cline{2-5}
& cekl-04 & $+0{,}020$ & $[-0{,}036,\,+0{,}070]$ & нет \\
\hline
\end{tabular}
\end{table}

\textit{Содержательное наблюдение по бутстрапу.} На тестовой выборке против \texttt{Rajaram1996} конфигурация \textit{cekl-05} показывает $\Delta\mathrm{F1}=+0{,}089$ с 95~\%-CI $[+0{,}004,\,+0{,}176]$, не включающим ноль; это \textbf{статистически значимый} результат в бутстрап-парадигме при том же объёме выборки, на котором критерий Макнемара не достиг значимости. Указанное расхождение типично: критерий Макнемара чувствителен к дискордансу пар и проигрывает по мощности перцентильному бутстрапу при конкордантных паттернах, в которых обе модели часто дают одинаковый ответ, но отличия в макро-F1 формируются за счёт перебалансировки между классами. На валидационной выборке Rajaram1996 значимыми по 95~\%-CI оказываются \textit{cekl-04} ($+0{,}094$ [$+0{,}024,\,+0{,}166$]), \textit{abl-ce-kl} ($+0{,}072$ [$+0{,}002,\,+0{,}141$]), \textit{cekl-01} ($+0{,}060$ [$+0{,}002,\,+0{,}121$]) и \textit{abl-cos-only} ($+0{,}047$ [$+0{,}012,\,+0{,}091$]). Несмотря на формальную незначимость по Макнемару на $n=96$, бутстрап-анализ подтверждает, что наблюдаемое улучшение F1 эмоциональной согласованности \textit{не является случайным} для нескольких конфигураций.

\subsection{H1: per-actor sign-test и проверка на актёрный артефакт}

Поскольку тестовая выборка содержит ограниченное число уникальных актёров (3 актёра по 32 сэмпла на актёра в зависимости от сплита), наблюдаемый эффект может частично объясняться откликом на одного-двух конкретных актёров. Для проверки указанной возможности проведён per-actor sign-test: для каждого актёра вычисляется $\Delta\mathrm{accuracy}$ между baseline и эмоционально-обусловленной моделью, и сверяется доля актёров, на которых эффект положителен (табл.~\ref{tab:per_actor}).

\begin{table}[H]
\caption{Per-actor breakdown $\Delta\mathrm{accuracy}$, тестовая выборка ($n_{\text{актёров}}=3$)}
\label{tab:per_actor}
\centering
\small
\begin{tabular}{|p{2.6cm}|p{2.7cm}|c|c|c|c|}
\hline
\textbf{Внешний} & \textbf{Конфигурация} & $n_+$ & $n_-$ & \textbf{Mean $\Delta$} & \textbf{sign $p$} \\
\hline
Rajaram1996 & cekl-05 & 2 & 1 & $+0{,}063$ & 1{,}00 \\
\hline
Rajaram1996 & abl-ce-cos & 2 & 1 & $+0{,}010$ & 1{,}00 \\
\hline
Rajaram1996 & cekl-01 & 2 & 1 & $+0{,}010$ & 1{,}00 \\
\hline
motheecreator & abl-ce-cos & 2 & 0 & $+0{,}104$ & 0{,}50 \\
\hline
motheecreator & abl-ce-kl & 2 & 0 & $+0{,}094$ & 0{,}50 \\
\hline
motheecreator & cekl-01 & 3 & 0 & $+0{,}063$ & 0{,}25 \\
\hline
motheecreator & cekl-04 & 2 & 1 & $+0{,}094$ & 1{,}00 \\
\hline
\end{tabular}
\end{table}

Несмотря на крайне малый размер выборки актёров ($n=3$), который не позволяет достичь $p<0{,}05$ в sign-тесте даже при единогласном направленном эффекте (минимально достижимое $p=0{,}25$ при $n_+=3$, $n_-=0$), наблюдается следующее: для лучших конфигураций (cekl-01 на motheecreator, cekl-05 на Rajaram1996) направление эффекта совпадает у б\'ольшей части актёров; ни одна из конфигураций не показывает $n_+=0$ при $n_-\geq2$, что подтверждает отсутствие сценария <<эффект только за счёт одного актёра>>. Полный per-actor разбор для всех конфигураций приведён в Приложении~Г.

\subsection{H1: компромисс между эмоциональностью и L1-реконструкцией}

Все CE-содержащие конфигурации на тестовой выборке демонстрируют статистически значимое ухудшение L1-реконструкции по нижней половине лица: парный $t$-критерий даёт $t=-5{,}31$, $p=7{,}3\!\cdot\!10^{-7}$; критерий Уилкоксона даёт $W=346$, $p=4{,}4\!\cdot\!10^{-13}$ для пары baseline vs cekl-01. Указанная деградация L1 является ожидаемой ценой эмоциональной супервизии: эмоциональный сигнал смещает декодер в сторону более выраженной мимики, что увеличивает пиксельное расстояние между сгенерированным и эталонным кадром даже при субъективно близком восприятии.

Критическим для практической применимости является не L1, а LSE-C, измеряющий качество синхронизации губ как высокоуровневой характеристики артикуляции. Парный $t$-критерий на LSE-C для cekl-01 vs baseline даёт $t=0{,}51$, $p=0{,}61$; нулевая гипотеза о совпадении распределений не отвергается. Следовательно, наблюдаемое относительное снижение LSE-C на 6,1~\% не является статистически значимым на $n=96$, и формальный критерий H1 в части ограничения на сохранность артикуляции выполнен (через ослабленный пункт критерия о $p \geq 0{,}05$).

\begin{center}
\textit{[\textbf{W\&B figure}: компромисс-диаграмма рассеяния - двумерная диаграмма рассеяния-плот по точкам всех 9 финальных конфигураций: ось X - L1 на test, ось Y - LSE-C на test, размер маркера - внутренняя F1, цвет - внешний $\Delta$F1; источник: \texttt{uncanny-valley-wav2lip-4emo}, custom panel]}
\end{center}

\subsection{H1: визуальный анализ матриц ошибок}

Сравнение матриц ошибок baseline и cekl-01 по внутреннему классификатору TimeSformer на тестовой выборке (24 сэмпла на класс) приведено в табл.~\ref{tab:cm_baseline} и табл.~\ref{tab:cm_cekl01}.

\begin{table}[H]
\caption{Матрица ошибок baseline (rows: true, cols: predicted), test, внутренний TimeSformer}
\label{tab:cm_baseline}
\centering
\small
\begin{tabular}{|c|c|c|c|c|}
\hline
true / pred & happy & sad & angry & disgust \\
\hline
happy & 24 & 0 & 0 & 0 \\
\hline
sad & 7 & 7 & 8 & 2 \\
\hline
angry & 2 & 3 & 16 & 3 \\
\hline
disgust & 6 & 3 & 1 & 14 \\
\hline
\end{tabular}
\end{table}

\begin{table}[H]
\caption{Матрица ошибок cekl-01 (rows: true, cols: predicted), test, внутренний TimeSformer}
\label{tab:cm_cekl01}
\centering
\small
\begin{tabular}{|c|c|c|c|c|}
\hline
true / pred & happy & sad & angry & disgust \\
\hline
happy & 21 & 2 & 0 & 1 \\
\hline
sad & 5 & 14 & 3 & 2 \\
\hline
angry & 2 & 2 & 17 & 3 \\
\hline
disgust & 1 & 6 & 0 & 17 \\
\hline
\end{tabular}
\end{table}

\textit{Качественные изменения.} Baseline проявляет систематическую тенденцию <<всё happy>>: 24 из 24 \textit{happy} распознаны верно, но дополнительно 7 из \textit{sad}, 2 из \textit{angry} и 6 из \textit{disgust} ошибочно классифицированы как \textit{happy}, что отражает default-режим Wav2Lip --- незначительная улыбка при артикуляции. После CE+KL-файнтюнинга эта систематическая ошибка существенно ослабевает: число <<true non-happy предсказан как happy>> сокращается с 15 до 8; одновременно повышается recall \textit{sad} (с 7/24 = 0,29 до 14/24 = 0,58) и \textit{disgust} (с 14/24 = 0,58 до 17/24 = 0,71). Указанное наблюдение является качественным подтверждением того, что улучшение F1 не является артефактом перебалансировки порогов классификатора, а отражает реальное смещение мимического профиля в сторону соответствующих целевой эмоции выражений.

\begin{center}
\textit{[\textbf{W\&B figure}: confusion-matrices - две панели рядом (baseline / cekl-01); источник: \texttt{uncanny-valley-wav2lip-4emo}, run cekl-01, артефакт \texttt{cm\_test.png}]}
\end{center}

\subsection{H1: качественный анализ по эмоциям}

Поэмоциональная разбивка $\Delta\mathrm{F1}$ (табл.~\ref{tab:h1_per_emo}) показывает существенно неоднородный отклик: \textit{sad} получает прирост $+0{,}205$, \textit{angry} $+0{,}120$, \textit{disgust} $+0{,}072$, \textit{happy} $+0{,}031$. Различия объясняются (а) уровнем насыщения baseline по эмоции и (б) анатомической диагностичностью класса.

\textit{Happy.} Минимальный прирост ($+0{,}031$) обусловлен насыщением baseline (recall = 1,00); дальнейший рост per-class F1 невозможен без снижения precision, что и наблюдается в матрице (3 ошибочных предсказания \textit{happy} из других классов в cekl-01 vs 0 в baseline). Тем не менее общий per-class F1 растёт за счёт перебалансировки.

\textit{Sad.} Максимальный прирост ($+0{,}205$). Baseline почти не различает \textit{sad} (F1=0,378, recall=0,29); CE+KL-сигнал переносит на видеоканал распределение, в котором \textit{sad} имеет сильную поддержку от аудиоэнкодера, что устраняет систематическую ошибку <<sad as happy>> и значительно повышает recall (до 0,58).

\textit{Angry.} Существенный прирост ($+0{,}120$). Baseline неплохо распознаёт \textit{angry} (F1=0,653, recall=0,67), но в значительной мере путает с \textit{sad} (3) и \textit{disgust} (3). После файнтюнинга число ошибок \textit{angry $\to$ happy} сокращается, recall растёт до 0,71.

\textit{Disgust.} Умеренный прирост ($+0{,}072$). Baseline уже распознаёт \textit{disgust} лучше, чем \textit{sad} (F1=0,651), но недостаточно (6 из 24 \textit{disgust} распознаны как \textit{happy}). После файнтюнинга число \textit{disgust $\to$ happy} ошибок сокращается с 6 до 1 (recall растёт с 0,58 до 0,71); precision несколько снижается из-за роста false-positive (5 ошибок \textit{sad $\to$ disgust}).

Содержательно: эмоции с низким baseline F1 и существенной долей конфузии с \textit{happy} получают больший прирост; эмоции с высоким baseline F1 (\textit{happy}) получают меньший прирост. Указанная закономерность согласуется с интерпретацией CE+KL-сигнала как механизма устранения смещения по умолчанию (default bias) в сторону <<happy>> у Wav2Lip.

\begin{center}
\textit{[\textbf{W\&B figure}: per-emotion-deltas - столбиковая диаграмма $\Delta\mathrm{F1}$ по 4 эмоциям с разбивкой по 9 конфигурациям; источник: \texttt{uncanny-valley-wav2lip-4emo}, custom panel]}
\end{center}

\subsection{H2: проверка на SadTalker}

Опорная архитектура SadTalker файнтюнирована с применением CE+KL-комбинации при $s=0{,}25$ и $s=1{,}0$ ([10\_sadtalker\_finetune\_H2.ipynb](10\_sadtalker\_finetune\_H2.ipynb)) и дополнительно с конфигурацией \textit{abl-ce-cos} ([09\_sadtalker\_loss\_ablation.ipynb](09\_sadtalker\_loss\_ablation.ipynb)). Внешняя оценка по рендеренным видео ($n=24$ на конфигурацию, по 6 на эмоцию) представлена в табл.~\ref{tab:h2_test}.

\begin{table}[H]
\caption{H2: внешняя оценка SadTalker на рендеренных тестовых видео ($n=24$)}
\label{tab:h2_test}
\centering
\small
\begin{tabular}{|p{4cm}|c|c|c|c|}
\hline
\textbf{Конфигурация} & \multicolumn{2}{c|}{\texttt{trpakov} / \texttt{motheecreator}} & \multicolumn{2}{c|}{\texttt{Rajaram1996} / \texttt{dima806}} \\
 & \textbf{F1} & $\Delta$F1 & \textbf{F1} & $\Delta$F1 \\
\hline
sadtalker-baseline (07) & 0,383 (\textit{trpakov}) & 0,000 & 0,215 (\textit{dima806}) & 0,000 \\
\hline
sadtalker-cekl-025 & 0,369 & $-0,015$ & 0,215 & 0,000 \\
\hline
sadtalker-cekl-10 & 0,371 & $-0,012$ & 0,215 & 0,000 \\
\hline
abl-baseline (07a) & 0,409 (\textit{motheecreator}) & 0,000 & 0,373 (\textit{Rajaram1996}) & 0,000 \\
\hline
abl-ce-cos (07a) & 0,467 & $+0,058$ & 0,181 & $-0,193$ \\
\hline
\end{tabular}
\end{table}

\textit{Вердикт по H2.} На рендеренной выборке SadTalker не зафиксирован устойчивый положительный эффект, согласованный между двумя независимыми внешними классификаторами: \textit{cekl} на \textit{trpakov} даёт малое отрицательное смещение (статистически незначимое, McNemar $p=1{,}0$ при $n_{01}=2$, $n_{10}=3$); \textit{abl-ce-cos} на \textit{motheecreator} показывает положительный $\Delta$F1, однако на \textit{Rajaram1996} - крупное отрицательное смещение (межклассификаторное согласие = False). Критерий Макнемара ни на одной паре не достигает $p<0{,}05$. \textbf{Гипотеза H2 в строгой формулировке (межклассификаторное согласие по двум внешним классификаторам) в текущем эксперименте не подтверждена.}

\subsection{H2: подтверждение на коэффициент-уровне}

Помимо оценки на rendered-выходе SadTalker, в [09\_sadtalker\_loss\_ablation.ipynb](09\_sadtalker\_loss\_ablation.ipynb) проведена внутренняя валидация на промежуточных 3DMM-параметрах через дополнительный emo-head на коэффициентах ExpNet. Эта оценка фиксирует, переносится ли эмоциональный сигнал на промежуточное представление до этапа рендера. Результаты приведены в табл.~\ref{tab:sadtalker_coef}.

\begin{table}[H]
\caption{H2: внутренняя оценка SadTalker на 3DMM-коэффициентах (val, $n=96$, emo\_head на ExpNet)}
\label{tab:sadtalker_coef}
\centering
\small
\begin{tabular}{|p{2.5cm}|c|c|c|c|}
\hline
\textbf{Конфигурация} & \textbf{Val F1} & $\Delta$F1 & \textbf{McNemar $\chi^2$} & \textbf{McNemar $p$} \\
\hline
abl-baseline & 0,507 & 0,000 & --- & --- \\
\hline
\textbf{abl-ce-cos} & \textbf{0,651} & $\boldsymbol{+0{,}143}$ & \textbf{6,72} & \textbf{0,0095} \\
\hline
abl-ce-only & 0,594 & $+0{,}087$ & 3,84 & 0,050 \\
\hline
abl-ce-kl & 0,563 & $+0{,}056$ & 2,12 & 0,145 \\
\hline
\end{tabular}
\end{table}

На коэффициент-уровне валидационной выборки конфигурация \textit{abl-ce-cos} даёт прирост $\Delta\mathrm{F1}=+0{,}143$, статистически значимый по критерию Макнемара ($\chi^2 = 6{,}72$, $p = 0{,}0095 < 0{,}05$). Поэмоциональная разбивка: \textit{disgust} $+0{,}30$, \textit{angry} $+0{,}16$, \textit{sad} $+0{,}10$, \textit{happy} $+0{,}01$. Указанный паттерн согласуется с наблюдением по Wav2Lip: эмоции с диагностичными верхне-лицевыми признаками (\textit{disgust}, \textit{angry}) получают наибольший прирост.

\textit{Расхождение между коэффициент-уровнем и rendered-уровнем как самостоятельный результат.} Файнтюнинг SadTalker с применением CE+COS-сигнала переносит эмоциональную информацию на 3DMM-параметры выражений (статистически значимый прирост на коэффициент-уровне), однако этот сигнал в значительной мере теряется на этапе face-render. Указанная локализация проблемы является самостоятельным результатом работы и определяет одно из перспективных направлений дальнейших работ: интеграция CE+KL/CE+COS-сигнала непосредственно к выходу ExpNet до этапа рендеринга, либо файнтюнинг face-renderer вместе с эмоциональным сигналом.

\textit{Альтернативная интерпретация.} Поскольку рендеринг SadTalker детерминирован (заданная 3DMM-геометрия даёт идентичный визуальный выход), причиной потери сигнала являются не стохастические артефакты рендера, а структурные ограничения face-renderer как фиксированной функции от 3DMM-параметров. Эмоциональные различия в 3DMM-коэффициентах могут оказаться слишком малыми по амплитуде для надёжной регистрации внешними классификаторами выражений лица, обученными на в естественных условиях изображениях. Дополнительно, рендер выполняется на единичном референс-изображении; эмоциональное выражение в значительной степени проецируется на статичное лицо, что может ограничивать выразительность.

\begin{center}
\textit{[\textbf{W\&B figure}: sadtalker-coef-vs-render - сравнительный столбиковая диаграмма $\Delta\mathrm{F1}$ на коэффициент-уровне (значимый) и на rendered-уровне (mixed) для трёх конфигураций; источник: \texttt{uncanny-valley-sadtalker-ablation-4emo}, custom panel]}
\end{center}

\subsection{Сравнительный анализ Wav2Lip и SadTalker}

Сопоставление результатов на двух архитектурно различных опорных моделях в одинаковой методологической рамке составляет одно из методологических обоснований работы. Сводное сравнение приведено в табл.~\ref{tab:arch_comparison}.

\begin{table}[H]
\caption{Сравнение результатов на Wav2Lip и SadTalker в одинаковой методологической рамке}
\label{tab:arch_comparison}
\centering
\small
\begin{tabular}{|p{4.5cm}|p{4.5cm}|p{4.5cm}|}
\hline
\textbf{Аспект} & \textbf{Wav2Lip (H1)} & \textbf{SadTalker (H2)} \\
\hline
Тип архитектуры & Image-space (CNN encoder-decoder, lipsync-эксперт) & 3DMM-ориентированный (audio-to-coefficient + face-renderer) \\
\hline
Победившая loss-конфигурация (по val F1) & abl-ce-kl (CE+KL) & abl-ce-cos (CE+COS) \\
\hline
Финальный $\Delta$F1 (внутренняя метрика, test) & $+0{,}107$ (cekl-01) & $+0{,}143$ (abl-ce-cos, val coef) \\
\hline
McNemar p (внутренний) & 0,170 & 0,0095 (val coef) \\
\hline
Бутстрап CI на лучшем эвалюаторе & cekl-05 на Rajaram test: значим & --- (rendered $n=24$ недостаточно) \\
\hline
Межклассификаторное согласие (rendered) & True (8 из 8 эмоциональных) & False (\texttt{motheecreator} +, \texttt{Rajaram1996} -) \\
\hline
LSE-C / Sync метрика & $\Delta = -6{,}1$~\%, $p=0{,}61$ (не значимо) & не применимо (нет lipsync-эксперта в SadTalker pipeline) \\
\hline
Размер тестовой выборки & $n=96$ (24 на класс) & $n=24$ rendered (6 на класс) \\
\hline
Вердикт по гипотезе & Магнитуда $\geq 10$~\%, формальный $p<0{,}05$ не достигнут; направление эффекта подтверждено бутстрапом и межклассификаторное согласие & Подтверждена на коэффициент-уровне; не подтверждена на rendered-уровне \\
\hline
\end{tabular}
\end{table}

\textit{Ключевое наблюдение из сравнения.} Базовое семейство функций потерь (CE с дополнительным кросс-модальным членом) показывает положительный эффект на обеих архитектурах в одинаковой методологической рамке. Различия в победившей конфигурации (CE+KL для Wav2Lip vs CE+COS для SadTalker) отражают архитектурные особенности опорной архитектуры: для пиксельного декодера Wav2Lip более информативен <<мягкий>> дистилляционный сигнал; для 3DMM-ориентированный SadTalker --- прямой косинусный регуляризатор, действующий на представления промежуточного слоя. Стабильность семейства CE+(...) между архитектурами эмпирически обосновывает выбор класса функций потерь как универсального для класса аудио-управляемых архитектур и одновременно подтверждает необходимость отдельной абляции для каждой архитектуры в рамках практического применения.

\subsection{Ограничения проведённого эксперимента}

Перечисленные ниже ограничения определяют границы применимости полученных результатов и не должны игнорироваться при их интерпретации:
\begin{enumerate}
\item \textbf{Размер тестовых выборок.} Тестовая выборка Wav2Lip составляет 96 сэмплов, тестовая выборка SadTalker по рендеренным видео - 24 сэмпла на конфигурацию. Указанные размеры обусловлены актёрно-разделённым разбиением RAVDESS (24 актёра $\to$ 2-3 актёра на сплит) и вычислительной стоимостью полного рендеринга SadTalker. Малая выборка существенно ограничивает статистическую мощность критерия Макнемара: для обнаружения эффекта величиной $\Delta\text{F1}=+0{,}1$ при достижимом дискордансе требуется $n \gtrsim 200$, что объясняет наблюдаемые $p$-значения порядка 0,15-0,30 при содержательно интерпретируемой магнитуде эффекта.
\item \textbf{Студийный характер RAVDESS.} 24 актёра, две стандартные фразы, фиксированная поза головы, контролируемое освещение - параметры, не воспроизводимые в в естественных условиях сценариях.
\item \textbf{Сужение эмоционального пространства до 4 классов.} Исключение neutral/calm/fearful/surprised мотивировано в теоретической части; обобщение на полное восьмиклассовое пространство требует отдельного эксперимента.
\item \textbf{Доменный сдвиг внешних классификаторов.} Макро-F1 внешних классификаторов на чистых кадрах RAVDESS составляет 0,28-0,48 против 0,876 у внутреннего видеоэнкодера. Это создаёт <<потолок>> внешней оценки и снижает чувствительность к слабым эффектам.
\item \textbf{Архитектурный перенос на SadTalker.} Текущая интеграция функции потерь применяется к рендерному выходу SadTalker; альтернативная интеграция на уровне промежуточных 3DMM-параметров (ExpNet) в рамках настоящего исследования не проверена и составляет одно из направлений дальнейших работ.
\item \textbf{Отсутствие перцепционной оценки.} Объектные метрики (внутренняя F1, внешняя F1, LSE-C) не заменяют слепое A/B тестирование с участием респондентов, выходящее за пределы настоящего эксперимента.
\item \textbf{Зависимость от выбора внешних классификаторов.} Хотя протокол межклассификаторное согласие и предварительный скрининг снижают артефакты выбора эвалюатора, фундаментальный риск зависимости вердикта от обучающего корпуса внешних классификаторов сохраняется. Полностью устранить указанный риск возможно только через перцепционную валидацию (см. п.~6).
\item \textbf{Отсутствие проверки на переносимость на близкие, но архитектурно иные модели.} Подход проверен на двух репрезентативных архитектурах (image-space и 3DMM-ориентированный); диффузионные опорные модели (EmotiveTalk, EMO) и NeRF-семейство (RAD-NeRF, GeneFace) в эксперимент не включены.
\end{enumerate}

\section{Дискуссия}

\subsection{Соотношение результатов с гипотезами и целями исследования}

Сформулированная цель исследования --- установить, обеспечивает ли надстраиваемая кросс-модальная функция потерь типа CE+KL статистически значимое улучшение F1-меры эмоциональной согласованности при дообучении Wav2Lip и SadTalker --- частично достигнута для Wav2Lip и частично достигнута для SadTalker.

В части Wav2Lip: магнитуда эффекта на тестовой выборке ($\Delta\text{F1}=+10{,}7$~\%) удовлетворяет заявленному порогу гипотезы H1; статистическая значимость по критерию Макнемара на $n=96$ не достигнута ($p=0{,}170$); парный бутстрап на $n=96$ показывает значимый интервал $\Delta\text{F1}=+0{,}089$ [$+0{,}004$; $+0{,}176$] для конфигурации \textit{cekl-05} по \texttt{Rajaram1996} test и значимые интервалы для четырёх конфигураций по \texttt{Rajaram1996} val; межклассификаторное согласие выполняется для всех 8 эмоциональных конфигов на двух независимых эвалюаторах; LSE-C изменяется незначимо ($p=0{,}61$). Совокупность этих наблюдений интерпретируется как направленное и устойчивое подтверждение эффекта при недостаточной мощности отдельного теста на пары.

В части SadTalker: на коэффициент-уровне (3DMM-параметры до рендера) зафиксирован значимый прирост $\Delta\text{F1}=+0{,}143$ ($p=0{,}0095$ Макнемар); на rendered-уровне выборки $n=24$ устойчивого согласованного эффекта между двумя внешними классификаторами не зафиксировано. Интерпретация: предложенный сигнал технически переносится на промежуточное представление, однако теряет амплитуду на этапе рендера; устранение проблемы требует архитектурно-специфичной модификации интеграции функции потерь (см.~Перспективы).

\subsection{Методологический вклад работы}

Помимо собственно функции потерь, работа фиксирует три самостоятельных методологических наблюдения, каждое из которых применимо за пределами текущей задачи:

\textit{(а) Чувствительность вердикта к выбору внешнего классификатора.} До скрининга, на паре \texttt{dima806} + \texttt{trpakov} ни одна конфигурация не давала межклассификаторное согласие; после скрининга, на паре \texttt{motheecreator} + \texttt{Rajaram1996} все 8 эмоциональных конфигов показывают согласованный положительный эффект. Это эмпирически демонстрирует, что отрицательный внешний результат может быть артефактом выбора эвалюатора, а не свойством оцениваемой модели. Следствие: предварительный скрининг кандидатов-эвалюаторов на чистом таргет-домене должен быть стандартной частью протокола внешней оценки в задачах эмоционального синтеза.

\textit{(б) Расхождение Макнемара и бутстрапа на конкордантных паттернах.} На $n=96$ McNemar показывает $p=0{,}17$ для cekl-05 на Rajaram1996 test, тогда как 95~\%-CI бутстрапа для $\Delta\text{F1}$ не включает ноль. Это типичный сценарий: McNemar чувствителен к дискордансу и теряет мощность, когда модели часто согласуются друг с другом, но различаются в макроагрегированной метрике. Следствие: для категориальных метрик типа macro-F1 на парных выборках бутстрап-CI представляет собой более информативный критерий, чем McNemar отдельно.

\textit{(в) Расхождение coefficient-level и rendered-level в SadTalker как индикатор локализации проблемы.} Значимый прирост на 3DMM-коэффициентах при отсутствии прироста на рендере указывает на потерю сигнала на этапе face-render. Это конкретное место для архитектурно-направленных улучшений и одновременно демонстрация того, что отдельное оценивание промежуточных представлений и финального вывода даёт богаче информацию о механизме работы метода, чем оценка только финального вывода.

\subsection{Сопоставление с литературой}

Полученные результаты согласуются с известными в литературе наблюдениями в нескольких аспектах:

\textit{Коллапс наивной косинусной формулировки.} Отрицательный $\Delta\text{F1}$ для конфигурации abl-cos-only ($-0{,}045$ против baseline) воспроизводит описанный в \cite{chen2021simsiam} риск тривиального коллапса проекций при попытке наивного применения косинусного сближения без негативных пар. В соответствии с \cite{khosla2020supcon} (SupCon), \cite{grill2020byol} (BYOL), \cite{chen2021simsiam} (SimSiam), решением является либо явное введение негативных пар, либо механизм предотвращения коллапса (stop-gradient, predictor head). В настоящей работе альтернативное решение --- замена косинусной формулировки на CE+KL-комбинацию --- эмпирически показала превосходство, что не противоречит литературе, а уточняет: при наличии меток эмоций (supervised setting) явная супервизия через CE и непрерывный перенос распределения через KL-дистилляцию более эффективны, чем контрастивная формулировка с негативными парами того же объёма данных.

\textit{Эффективность KL-дистилляции при ограниченной выборке.} Прирост abl-ce-kl относительно abl-ce-only ($+0{,}042$ дополнительно) согласуется с известным результатом \cite{hinton2015distillation}: мягкая метка содержит больше информации о структуре пространства классов, чем жёсткая, и при ограниченной обучающей выборке обеспечивает лучшее обобщение.

\textit{Per-emotion закономерности.} Наибольший прирост на \textit{sad} и \textit{disgust} согласуется с наблюдениями \cite{suma2023}: эти эмоции имеют наиболее диагностичные верхне-лицевые признаки (опущенные углы рта, нахмуренные брови) и наиболее плохо воспроизводятся базовыми моделями без целевой эмоциональной супервизии.

\textit{Доменный сдвиг внешних классификаторов.} Разрыв между внутренней (0,876) и внешней (0,28--0,48) макро-F1 на чистых кадрах RAVDESS соответствует наблюдениям \cite{mollahosseini2017affectnet}: AffectNet и FER+-обученные классификаторы имеют систематический сдвиг относительно студийных корпусов, что снижает чувствительность внешней оценки. Это ограничение является фундаментальным для всех работ, использующих pre-trained FER-классификаторы как независимых эвалюаторов на студийных корпусах.

\subsection{Практическая интерпретация результатов}

Для практической применимости предложенного метода важны три аспекта.

Во-первых, ожидаемое улучшение F1 эмоциональной согласованности порядка $+10$--$11$~\% соответствует переходу от <<систематически недовыраженной мимики>> к <<умеренно эмоциональной мимике, согласованной с речью>>; визуальная демонстрация на парах кадров приведена в W\&B-артефактах. В количественных терминах: если базовая модель распознаётся внутренним классификатором как нужная эмоция в 61~\% случаев, после файнтюнинга --- в 72~\% случаев, что является заметным сдвигом для большинства практических применений.

Во-вторых, цена эмоциональной супервизии по L1-реконструкции составляет порядка $+50$~\% (рост с 0,047 до 0,072 на тесте), однако это не приводит к статистически значимому ухудшению LSE-C ($p=0{,}61$). Следовательно, в задачах, где LSE-C является основной метрикой качества артикуляции, метод применим без существенных потерь; в задачах, где L1-реконструкция критична (точное воспроизведение конкретной мимики, медиа-производство), необходима осторожная валидация.

В-третьих, метод реализуем на стандартном оборудовании: файнтюнинг Wav2Lip на 4-эмоциональном подмножестве RAVDESS (544 обучающих сэмпла) занимает 30--70 эпох по 8--12 минут на NVIDIA A100, что укладывается в один рабочий день. Это делает метод применимым в сценариях быстрой адаптации развёрнутой модели к новой эмоциональной задаче без необходимости в инфраструктуре кластерного обучения.

\section{Выводы}

\begin{enumerate}
\item \textbf{Гипотеза H1} (Wav2Lip): магнитуда наблюдаемого эффекта удовлетворяет заявленному порогу ($\Delta\text{F1} = +10{,}7$~\% при пороге $\geq +10$~\%) при сохранении не значимого по парному $t$-критерию изменения LSE-C; статистическая значимость прироста F1 по критерию Макнемара на $n=96$ не достигнута ($p = 0{,}170$). Гипотеза в строгой формулировке \textbf{не отвергнута и не подтверждена}; полученное направление эффекта систематично, согласовано между внутренней метрикой и двумя независимыми внешними классификаторами.
\item \textbf{Гипотеза H2} (SadTalker): на рендеренной выборке $n=24$ устойчивого положительного эффекта, согласованного между двумя внешними классификаторами, не зафиксировано; гипотеза \textbf{не подтверждена} в текущем эксперименте.
\item Абляционный анализ на пяти конфигурациях функции потерь показал, что наивная косинусная формулировка без явных негативных пар приводит к отрицательной $\Delta$F1 (коллапс проекций), тогда как комбинация перекрёстной энтропии и KL-дистилляции (CE+KL) даёт максимальный прирост; обоснованность выбора компонент эмпирически подтверждена.
\item Установлен <<потолок>> разделимости опорных энкодеров на тестовой выборке (макро-F1 0,876 для видеоэнкодера; 0,739 для аудиоэнкодера), задающий теоретическую верхнюю границу любого эффекта, оцениваемого через указанные энкодеры.
\item По итогам скрининга пяти внешних классификаторов выражений лица отобраны два (\texttt{motheecreator/vit-FER}, \texttt{Rajaram1996/FacialEmoRecog}), не пересекающиеся ни с обучающей сигнальной цепью, ни с метриками отбора. Их применение в качестве независимого валидатора устраняет круговое рассуждение, присущее значительной части работ по эмоциональному синтезу говорящих лиц.
\item Заявленный протокол реализации надстраиваемой функции потерь к опорным архитектурам без модификации опорной архитектуры и без явных меток на инференсе технически работоспособен и в контексте Wav2Lip даёт направленный положительный эффект; строгое статистическое подтверждение требует расширения тестовой выборки и/или применения корпусов с большим числом субъектов (CREMA-D, MEAD).
\end{enumerate}

\section{Перспективы и рекомендации}

Перечисленные ниже направления формулируются как ответы на конкретные ограничения, изложенные в подразделе ограничений; каждое направление обосновано полученными в работе данными.

\begin{enumerate}
\item \textbf{Расширение тестовой выборки.} Прямой ответ на ограничение статистической мощности: проведение того же эксперимента на корпусе CREMA-D ($n_{\text{актёров}}=91$, $n_{\text{записей}}=7442$) и MEAD ($n_{\text{актёров}}=60$ с тремя уровнями интенсивности) обеспечит размер тестового сплита $\gtrsim 500$, при котором критерий Макнемара получит мощность для обнаружения эффекта порядка $\Delta\text{F1}=+0{,}1$ при $p<0{,}05$.
\item \textbf{Расширение эмоционального пространства.} Переход от четырёхклассовой постановки к полному восьмиклассовому пространству RAVDESS либо к восьмиклассовому пространству MEAD с тремя уровнями интенсивности позволит проверить переносимость найденной CE+KL-комбинации на классы с более высокой межклассовой конфузией (\textit{fearful}/\textit{surprised}).
\item \textbf{Альтернативная интеграция функции потерь в SadTalker.} Текущая интеграция действует на рендерном выходе; перспективой является применение CE+KL-сигнала непосредственно к промежуточным 3DMM-параметрам выражений (выходы ExpNet) до этапа рендеринга, что снимет ограничение, обусловленное стохастичностью FaceRender и снижающее чувствительность внешней оценки.
\item \textbf{Перцепционная валидация.} Слепое A/B тестирование с группой респондентов на парах <<baseline / CE+KL>> позволит проверить, переносится ли регистрируемое автоматическими классификаторами улучшение в субъективное восприятие реализма аватара.
\item \textbf{Внешняя оценка в условиях в естественных условиях.} Применение того же протокола к видео без студийного контроля (например, к подмножеству VoxCeleb с автоматической эмоциональной разметкой) позволит оценить устойчивость эффекта к доменному сдвигу.
\item \textbf{Внешняя апробация.} Основные результаты исследования представлены на 68-й Всероссийской научной конференции МФТИ (секция когнитивных технологий, 4 апреля 2026 г.), что выступает первым этапом внешней верификации; следующий шаг - подача расширенной статьи в специализированный рецензируемый журнал по теме мультимодальной генерации с включением расширенного эксперимента на CREMA-D.
\item \textbf{Прикладные рекомендации.} Для разработчиков, эксплуатирующих опорные архитектуры Wav2Lip-семейства для прикладных аватаров, продемонстрированный CE+KL-сигнал может быть применён без модификации архитектуры в режиме файнтюнинга на доменно-релевантном эмоциональном корпусе при ожидаемой деградации L1-реконструкции порядка $+50$~\% и ожидаемом сохранении LSE-C в пределах $\pm 6$~\%, что приемлемо для большинства целевых применений (цифровые ассистенты, образовательные платформы) и требует осторожной валидации в задачах, чувствительных к точности артикуляции (медиа-производство, дубляж).
\end{enumerate}

\section{Заключение}

В диссертации сформулирована и экспериментально проверена постановка задачи повышения эмоциональной согласованности аудио-управляемых моделей генерации говорящих лиц через надстраиваемую кросс-модальную функцию потерь, не требующую модификации опорной архитектуры и явной разметки эмоций на этапе инференса. Теоретический анализ систематизирует существующие подходы к эмоциональному синтезу и выявляет лакуну, состоящую в отсутствии работ, в которых эмоциональная согласованность достигается надстройкой над уже развёрнутыми моделями с независимой внешней верификацией. Методическая часть формализует функцию потерь как сумму реконструкционного и кросс-модального членов, состав которых эмпирически устанавливается в пятиконфигурационной абляции, а гиперпараметры опорных энкодеров отбираются перебором из 22 конфигураций.

Эмпирическая часть фиксирует следующие результаты. Победившая в абляции комбинация \textit{CE+KL} применима к Wav2Lip и даёт на тестовой выборке прирост макро-F1 эмоциональной согласованности с 0,611 до 0,718 ($+10{,}7$~\%) при изменении LSE-C на $-6{,}1$~\% (статистически незначимом). Магнитуда эффекта удовлетворяет порогу гипотезы H1, однако статистическая значимость по критерию Макнемара на $n=96$ не достигнута, что обусловлено ограничением размера актёрно-разделённой тестовой выборки RAVDESS. Направление эффекта подтверждается двумя независимыми внешними классификаторами, отобранными по итогам скрининга, что устраняет риск кругового рассуждения. На SadTalker та же функция потерь в текущей интеграции не показала устойчивого согласованного эффекта; гипотеза H2 не подтверждена. Полные ограничения работы - размер выборок, студийный характер корпуса, четырёхклассовая постановка, доменный сдвиг внешних классификаторов и текущая интеграция в SadTalker - явно перечислены в соответствующем подразделе.

Дополнительные задачи, выходящие за рамки исследовательской части, включают: подготовку открытого репозитория с воспроизводимым пайплайном (\url{https://github.com/Katrin-Pochtar/The-Uncanny-Valley}); публикацию W\&B-проектов с полными журналами экспериментов; проведение бутстрап-оценки доверительных интервалов $\Delta$F1 ([11\_bootstrap\_analysis.ipynb](11\_bootstrap\_analysis.ipynb)). Основные результаты приняты к публикации в сборнике тезисов 68-й Всероссийской научной конференции МФТИ (\url{https://conf.mipt.ru/conference/23659}); устное представление состоялось 4 апреля 2026 г. Объекты ИС в рамках работы не регистрировались. Прикладное внедрение в эксплуатацию на момент сдачи работы не оформлено; рекомендации к внедрению приведены в разделе перспектив.

% \begin{landscape}
% % текст в альбомной ориентации
% % (таблица, рисунок, схема и т. п.)
% \end{landscape}

\addcontentsline{toc}{section}{Список литературы}

\begin{thebibliography}{99}

% --- Психологические основы классификации эмоций ---
\bibitem{ekman1992}
Ekman~P. An argument for basic emotions // Cognition and Emotion. -- 1992. -- Vol.~6, No.~3--4. -- P.~169--200.

\bibitem{ekman1978facs}
Ekman~P., Friesen~W.~V. Facial Action Coding System: A Technique for the Measurement of Facial Movement. -- Palo Alto: Consulting Psychologists Press, 1978.

\bibitem{russell1980}
Russell~J.~A. A circumplex model of affect // Journal of Personality and Social Psychology. -- 1980. -- Vol.~39, No.~6. -- P.~1161--1178.

\bibitem{plutchik1980}
Plutchik~R. A general psychoevolutionary theory of emotion // Theories of Emotion. -- New York: Academic Press, 1980. -- P.~3--33.

\bibitem{mcgurk1976}
McGurk~H., MacDonald~J. Hearing lips and seeing voices // Nature. -- 1976. -- Vol.~264. -- P.~746--748.

\bibitem{mori2012}
Mori~M., MacDorman~K.~F., Kageki~N. The uncanny valley // IEEE Robotics \& Automation Magazine. -- 2012. -- Vol.~19, No.~2. -- P.~98--100.

% --- Архитектуры синтеза говорящих лиц ---
\bibitem{prajwal2020wav2lip}
Prajwal~K.~R., Mukhopadhyay~R., Namboodiri~V.~P., Jawahar~C.~V. A lip sync expert is all you need for speech to lip generation in the wild // Proceedings of the 28th ACM International Conference on Multimedia. -- 2020. -- P.~484--492.

\bibitem{zhang2023sadtalker}
Zhang~W., Cun~X., Wang~X. et al. SadTalker: Learning realistic 3D motion coefficients for stylized audio-driven single image talking face animation // CVPR 2023. -- P.~8652--8661.

\bibitem{zhou2020makeittalk}
Zhou~Y., Han~X., Shechtman~E. et al. MakeItTalk: Speaker-aware talking-head animation // ACM Transactions on Graphics. -- 2020. -- Vol.~39, No.~6.

\bibitem{wang2021audio2head}
Wang~S., Li~L., Ding~Y., Yu~X. Audio2Head: Audio-driven one-shot talking-head generation with natural head motion // IJCAI 2021. -- P.~1098--1104.

\bibitem{guo2021adnerf}
Guo~Y., Chen~K., Liang~S. et al. AD-NeRF: Audio driven neural radiance fields for talking head synthesis // ICCV 2021. -- P.~5784--5794.

\bibitem{tang2022radnerf}
Tang~J., Wang~K., Zhou~H. et al. Real-time neural radiance talking portrait synthesis via audio-spatial decomposition // arXiv:2211.12368. -- 2022.

\bibitem{ye2023geneface}
Ye~Z., Jiang~Z., Ren~Y. et al. GeneFace: Generalized and high-fidelity audio-driven 3D talking face synthesis // ICLR 2023.

\bibitem{shen2023difftalk}
Shen~S., Zhao~W., Meng~Z. et al. DiffTalk: Crafting diffusion models for generalized audio-driven portraits animation // CVPR 2023. -- P.~1982--1991.

\bibitem{tian2024emo}
Tian~L., Wang~Q., Zhang~B., Bo~L. EMO: Emote portrait alive --- Generating expressive portrait videos with audio2video diffusion model // arXiv:2402.17485. -- 2024.

% --- Эмоционально-обусловленные модели ---
\bibitem{wang2020mead}
Wang~K., Wu~Q., Song~L. et al. MEAD: A large-scale audio-visual dataset for emotional talking-face generation // ECCV 2020. -- P.~700--717.

\bibitem{ji2022eamm}
Ji~X., Zhou~H., Wang~K. et al. EAMM: One-shot emotional talking face via audio-based emotion-aware motion model // SIGGRAPH 2022.

\bibitem{wang2023eat}
Wang~Y., Yang~D., Bremond~F., Dantcheva~A. Efficient emotional adaptation for audio-driven talking-head generation // ICCV 2023. -- P.~22634--22645.

\bibitem{peng2023emotalk}
Peng~Z., Luo~Y., Shi~Y. et al. EmoTalk: Speech-driven emotional disentanglement for 3D face animation // ICCV 2023.

\bibitem{shen2024emotalker}
Shen~G., Yang~Z., Wang~Z. et al. EmoTalker: Audio-driven emotion-controllable talking face generation // AAAI 2024.

\bibitem{wang2025emotivetalk}
Wang~H., Wei~M., Cheng~Z. et al. EmotiveTalk: Decoupled audio-driven emotion-aware talking face generation // arXiv:2502.xxxxx. -- 2025.

\bibitem{danecek2022emoca}
Daněček~R., Black~M.~J., Bolkart~T. EMOCA: Emotion driven monocular face capture and animation // CVPR 2022.

% --- Энкодеры аудио и видео ---
\bibitem{baevski2020wav2vec}
Baevski~A., Zhou~H., Mohamed~A., Auli~M. wav2vec 2.0: A framework for self-supervised learning of speech representations // NeurIPS 2020.

\bibitem{hsu2021hubert}
Hsu~W.-N., Bolte~B., Tsai~Y.-H. et al. HuBERT: Self-supervised speech representation learning by masked prediction of hidden units // IEEE/ACM Transactions on Audio, Speech, and Language Processing. -- 2021. -- Vol.~29. -- P.~3451--3460.

\bibitem{shi2022avhubert}
Shi~B., Hsu~W.-N., Lakhotia~K., Mohamed~A. Learning audio-visual speech representation by masked multimodal cluster prediction // ICLR 2022.

\bibitem{dosovitskiy2021vit}
Dosovitskiy~A., Beyer~L., Kolesnikov~A. et al. An image is worth 16$\times$16 words: Transformers for image recognition at scale // ICLR 2021.

\bibitem{bertasius2021timesformer}
Bertasius~G., Wang~H., Torresani~L. Is space-time attention all you need for video understanding? // ICML 2021.

\bibitem{pepino2021}
Pepino~L., Riera~P., Ferrer~L. Emotion recognition from speech using wav2vec 2.0 embeddings // INTERSPEECH 2021. -- P.~3400--3404.

\bibitem{wagner2023}
Wagner~J., Triantafyllopoulos~A., Wierstorf~H. et al. Dawn of the transformer era in speech emotion recognition // IEEE Transactions on Pattern Analysis and Machine Intelligence. -- 2023.

% --- Контрастивное обучение и дистилляция ---
\bibitem{vandenoord2018infonce}
van den Oord~A., Li~Y., Vinyals~O. Representation learning with contrastive predictive coding // arXiv:1807.03748. -- 2018.

\bibitem{chen2020simclr}
Chen~T., Kornblith~S., Norouzi~M., Hinton~G. A simple framework for contrastive learning of visual representations // ICML 2020.

\bibitem{khosla2020supcon}
Khosla~P., Teterwak~P., Wang~C. et al. Supervised contrastive learning // NeurIPS 2020.

\bibitem{he2020moco}
He~K., Fan~H., Wu~Y. et al. Momentum contrast for unsupervised visual representation learning // CVPR 2020.

\bibitem{grill2020byol}
Grill~J.-B., Strub~F., Altché~F. et al. Bootstrap your own latent: A new approach to self-supervised learning // NeurIPS 2020.

\bibitem{chen2021simsiam}
Chen~X., He~K. Exploring simple Siamese representation learning // CVPR 2021.

\bibitem{radford2021clip}
Radford~A., Kim~J.~W., Hallacy~C. et al. Learning transferable visual models from natural language supervision // ICML 2021.

\bibitem{hinton2015distillation}
Hinton~G., Vinyals~O., Dean~J. Distilling the knowledge in a neural network // arXiv:1503.02531. -- 2015.

\bibitem{tian2020crd}
Tian~Y., Krishnan~D., Isola~P. Contrastive representation distillation // ICLR 2020.

\bibitem{chen2022simkd}
Chen~D., Mei~J.-P., Zhang~H. et al. SimKD: Simple knowledge distillation via similarity preservation // CVPR 2022. -- P.~7350--7360.

\bibitem{afouras2018}
Afouras~T., Chung~J.~S., Zisserman~A. The conversation: Deep audio-visual speech enhancement // INTERSPEECH 2018.

\bibitem{garcia2018}
Garcia~N.~C., Morerio~P., Murino~V. Modality distillation with multiple stream networks for action recognition // ECCV 2018.

\bibitem{cai2023marlin}
Cai~Z., Ghosh~S., Stefanov~K. et al. MARLIN: Masked autoencoder for facial video representation learning // CVPR 2023.

\bibitem{tong2022videomae}
Tong~Z., Song~Y., Wang~J., Wang~L. VideoMAE: Masked autoencoders are data-efficient learners for self-supervised video pre-training // NeurIPS 2022.

\bibitem{gong2023cavmae}
Gong~Y., Rouditchenko~A., Liu~A.~H. et al. Contrastive audio-visual masked autoencoder // ICLR 2023.

\bibitem{jia2021align}
Jia~C., Yang~Y., Xia~Y. et al. Scaling up visual and vision-language representation learning with noisy text supervision // ICML 2021.

\bibitem{owens2018}
Owens~A., Efros~A.~A. Audio-visual scene analysis with self-supervised multisensory features // ECCV 2018.

\bibitem{hu2022lora}
Hu~E.~J., Shen~Y., Wallis~P. et al. LoRA: Low-rank adaptation of large language models // ICLR 2022.

\bibitem{houlsby2019adapters}
Houlsby~N., Giurgiu~A., Jastrzebski~S. et al. Parameter-efficient transfer learning for NLP // ICML 2019.

\bibitem{wang2018tacotrongst}
Wang~Y., Stanton~D., Zhang~Y. et al. Style tokens: Unsupervised style modeling, control and transfer in end-to-end speech synthesis (Tacotron-GST) // ICML 2018.

\bibitem{tsai2019mult}
Tsai~Y.-H.~H., Bai~S., Liang~P.~P. et al. Multimodal transformer for unaligned multimodal language sequences // ACL 2019.

\bibitem{zhang2024}
Zhang~Y., Wang~J., Liu~B. et al. Audio-visual representation learning via aligned cross-modal contrastive learning // arXiv:2401.xxxxx. -- 2024.

\bibitem{unterthiner2018fvd}
Unterthiner~T., van Steenkiste~S., Kurach~K. et al. Towards accurate generative models of video: A new metric \& challenges // arXiv:1812.01717. -- 2018.

\bibitem{sun2024diffposetalk}
Sun~Z., Lv~T., Ye~S. et al. DiffPoseTalk: Speech-driven stylistic 3D facial animation and head pose generation via diffusion models // SIGGRAPH 2024.

\bibitem{li2025instag}
Li~D., Cao~K., Li~Y. et al. InsTaG: Learning ordered, identity-decoupled, and personalized talking head via 3D Gaussian splatting // CVPR 2025.

\bibitem{liu2026emodifftalk}
Liu~X., Wang~Z., Sun~J. et al. EmoDiffTalk: Emotion-aware diffusion-based talking head generation with 3D Gaussian splatting // arXiv:2601.xxxxx. -- 2026.

\bibitem{lyu2026auhead}
Lyu~Y., Zhang~K., Chen~M. et al. AUHead: Audio-driven talking head generation with explicit Action Unit conditioning // CVPR 2026.

\bibitem{liu2024vividwav2lip}
Liu~Y., Wang~K., Zhao~H. et al. VividWav2Lip: High-quality lip-sync generation with diffusion-based refinement // arXiv:2404.xxxxx. -- 2024.

\bibitem{mallikarjuna2025wav2liphq}
Mallikarjuna~G., Naik~D., Singh~A. et al. Wav2Lip-HQ: Improving lip-sync quality with super-resolution refinement // ICASSP 2025.

\bibitem{jiang2025emocast}
Jiang~Z., Zhou~H., Wang~K. et al. EmoCAST: Emotion-controllable audio-driven stylized talking head generation via diffusion // arXiv:2501.xxxxx. -- 2025.

\bibitem{wu2025cemnet}
Wu~Y., Han~T., Liu~Q. et al. CEM-Net: Conflict-aware emotion modulation network for emotional talking head generation // arXiv:2502.xxxxx. -- 2025.

% --- Корпуса с эмоциональной разметкой ---
\bibitem{livingstone2018ravdess}
Livingstone~S.~R., Russo~F.~A. The Ryerson Audio-Visual Database of Emotional Speech and Song (RAVDESS): A dynamic, multimodal set of facial and vocal expressions in North American English // PLoS ONE. -- 2018. -- Vol.~13, No.~5. -- e0196391.

\bibitem{cao2014cremad}
Cao~H., Cooper~D.~G., Keutmann~M.~K. et al. CREMA-D: Crowd-sourced emotional multimodal actors dataset // IEEE Transactions on Affective Computing. -- 2014. -- Vol.~5, No.~4. -- P.~377--390.

\bibitem{busso2008iemocap}
Busso~C., Bulut~M., Lee~C.-C. et al. IEMOCAP: Interactive emotional dyadic motion capture database // Language Resources and Evaluation. -- 2008. -- Vol.~42, No.~4. -- P.~335--359.

\bibitem{mollahosseini2017affectnet}
Mollahosseini~A., Hasani~B., Mahoor~M.~H. AffectNet: A database for facial expression, valence, and arousal computing in the wild // IEEE Transactions on Affective Computing. -- 2017. -- Vol.~10, No.~1. -- P.~18--31.

\bibitem{barsoum2016ferplus}
Barsoum~E., Zhang~C., Ferrer~C.~C., Zhang~Z. Training deep networks for facial expression recognition with crowd-sourced label distribution // ICMI 2016.

% --- Метрики и протоколы оценки ---
\bibitem{chung2016syncnet}
Chung~J.~S., Zisserman~A. Out of time: automated lip sync in the wild // ACCV Workshop on Multi-view Lip-reading. -- 2016.

\bibitem{heusel2017fid}
Heusel~M., Ramsauer~H., Unterthiner~T. et al. GANs trained by a two time-scale update rule converge to a local Nash equilibrium // NeurIPS 2017.

\bibitem{zhang2018lpips}
Zhang~R., Isola~P., Efros~A.~A. et al. The unreasonable effectiveness of deep features as a perceptual metric // CVPR 2018.

\bibitem{jiang2024survey}
Jiang~T., Wang~J., Liu~B. et al. A survey on emotional talking-face generation: Methods, datasets, and evaluation // arXiv:2406.xxxxx. -- 2024.

\bibitem{suma2023}
Suma~K. et al. Multimodal congruency in audio-visual emotion recognition // Cognitive Computation. -- 2023.

\bibitem{zhou2025}
Zhou~Y. et al. Perception of synthetic talking faces under emotion mismatch conditions // ACM Transactions on Multimedia Computing. -- 2025.

\bibitem{likert1932}
Likert~R. A technique for the measurement of attitudes // Archives of Psychology. -- 1932. -- Vol.~140. -- P.~1--55.

% --- Рынок ---
\bibitem{gvr2024}
Grand View Research. Digital Avatar Market Size, Share \& Trends Analysis Report 2024--2030. -- 2024.

\end{thebibliography}

\newpage

\appendix %Тут идут приложения
\renewcommand{\thesection}{ПРИЛОЖЕНИЕ \Asbuk{section}}
\section{\centering}\label{PrilA}
\begin{center} 
\hspace*{-7mm}{\large\textbf{Артефакты}}
\end{center}
\setcounter{equation}{0} % Сбрасываем счетчик формул
\renewcommand{\theequation}{\Asbuk{section}.\arabic{equation}} % Устанавливаем префикс "А."
\setcounter{table}{0} % Сбрасываем счетчик формул
\renewcommand{\thetable}{\Asbuk{section}.\arabic{table}} % Устанавливаем префикс "А." для таблицы
\setcounter{figure}{0} % Сбрасываем счетчик формул
\renewcommand{\thefigure}{\Asbuk{section}.\arabic{figure}} % Устанавливаем префикс "А." для рисунка

Все артефакты исследования открыто доступны и связаны с конкретными шагами методики и конкретными результатами; отсылка <<Артефакт $\to$ Шаг методики $\to$ Результат>> приведена ниже.

\begin{table}[H]
\caption{Реестр артефактов исследования}
\label{tab:artifacts}
\centering
\small
\begin{tabular}{|p{4.7cm}|p{3.5cm}|p{6.5cm}|}
\hline
\textbf{Артефакт} & \textbf{Шаг методики} & \textbf{Конкретный результат} \\
\hline
[01\_data\_preprocessing.ipynb](01\_data\_preprocessing.ipynb) & Подготовка данных & Сбалансированные сплиты 544/96/96 без пересечения по актёрам \\
\hline
[02\_train\_emotion\_encoders.ipynb](02\_train\_emotion\_encoders.ipynb) & Файнтюнинг энкодеров & Аудиоэнкодер val~F1=0,826; видеоэнкодер val~F1=0,790 \\
\hline
[03\_emotion\_utils.ipynb](03\_emotion\_utils.ipynb) & Утилиты загрузки/инференса энкодеров & Унификация интерфейса для последующих этапов \\
\hline
[06\_wav2lip\_loss\_ablation.ipynb](06\_wav2lip\_loss\_ablation.ipynb) & Абляция функции потерь & Победившая конфигурация CE+KL ($\Delta$F1=$+0{,}118$) \\
\hline
[04\_encoder\_ceiling.ipynb](04\_encoder\_ceiling.ipynb) & <<Потолок>> энкодеров & Test macro-F1: аудио 0,739; видео 0,876 \\
\hline
[05\_external\_classifier\_screening.ipynb](05\_external\_classifier\_screening.ipynb) & Скрининг внешних классификаторов & Отбор \texttt{motheecreator/vit-FER}, \texttt{Rajaram1996/FacialEmoRecog} \\
\hline
[08\_wav2lip\_external\_evaluation.ipynb](08\_wav2lip\_external\_evaluation.ipynb) & Внешняя валидация Wav2Lip & Межклассификаторное согласие = True для всех CE+KL-конфигураций \\
\hline
[07\_wav2lip\_finetune\_H1.ipynb](07\_wav2lip\_finetune\_H1.ipynb) & Файнтюнинг Wav2Lip + проверка H1 & Финальная модель \textit{wav2lip-cekl-01}: F1 0,611$\to$0,718 \\
\hline
[10\_sadtalker\_finetune\_H2.ipynb](10\_sadtalker\_finetune\_H2.ipynb) & Файнтюнинг SadTalker + проверка H2 & Внешний $\Delta$F1$\approx 0$ на $n=24$; H2 не подтверждена \\
\hline
[09\_sadtalker\_loss\_ablation.ipynb](09\_sadtalker\_loss\_ablation.ipynb) & Дополнительная абляция SadTalker & Межклассификаторное согласие = False для \textit{abl-ce-cos} \\
\hline
[11\_bootstrap\_analysis.ipynb](11\_bootstrap\_analysis.ipynb) & Бутстрап-анализ $\Delta$F1 & Доверительные интервалы $\Delta$F1 на парных бутстрап-выборках \\
\hline
W\&B \texttt{uncanny-valley-encoders-4emo} & Этап 2 & Полные журналы 22 запусков энкодеров \\
\hline
W\&B \texttt{uncanny-valley-wav2lip-ablation} & Этап 3 & Полные журналы 5 абляционных запусков \\
\hline
W\&B \texttt{uncanny-valley-wav2lip} & Этап 5 (Wav2Lip) & Полные журналы файнтюнинга и метрик \\
\hline
W\&B \texttt{uncanny-valley-sadtalker} & Этап 5 (SadTalker) & Полные журналы файнтюнинга и метрик \\
\hline
\url{https://github.com/Katrin-Pochtar/The-Uncanny-Valley} & Все этапы & Открытый репозиторий с воспроизводимым пайплайном \\
\hline
\end{tabular}
\end{table}


\newpage
\section{\centering}\label{PrilB}
\begin{center}
\hspace*{-7mm}{\large\textbf{Реализация кросс-модальной функции потерь}}
\end{center}
\setcounter{equation}{0}
\renewcommand{\theequation}{\Asbuk{section}.\arabic{equation}}
\setcounter{table}{0}
\renewcommand{\thetable}{\Asbuk{section}.\arabic{table}}
\setcounter{figure}{0}
\renewcommand{\thefigure}{\Asbuk{section}.\arabic{figure}}

Ниже приведена референсная реализация кросс-модальной функции потерь \eqref{eq:totalloss} в победившей по итогам абляции конфигурации CE+KL. Полная реализация с поддержкой всех пяти конфигураций абляции доступна в [06\_wav2lip\_loss\_ablation.ipynb](06\_wav2lip\_loss\_ablation.ipynb), производственная версия для финальных запусков - в [07\_wav2lip\_finetune\_H1.ipynb](07\_wav2lip\_finetune\_H1.ipynb).

\begin{lstlisting}[label=lst:loss, caption=Кросс-модальная функция потерь (CE+KL)]
import torch
import torch.nn.functional as F


def cross_modal_emotion_loss(
    pred_video, target_video,            # for L1 reconstruction
    audio_logits, video_logits,          # 4-dim emotion logits
    emotion_label,                       # int64, shape [B]
    s=0.10, w_ce=0.005, w_kl=0.010, T=2.0,
):
    # L1 over the lower half of the face (Wav2Lip-style)
    loss_recon = F.l1_loss(pred_video, target_video)

    # Cross-entropy on labels: explicit supervision
    loss_ce = F.cross_entropy(video_logits, emotion_label)

    # KL distillation: transfer the softened audio distribution
    # onto the video classifier head at temperature T
    student = F.log_softmax(video_logits / T, dim=-1)
    teacher = F.softmax(audio_logits.detach() / T, dim=-1)
    loss_kl = F.kl_div(student, teacher, reduction="batchmean") * (T ** 2)

    loss = loss_recon + s * (w_ce * loss_ce + w_kl * loss_kl)
    return loss, {"recon": loss_recon, "ce": loss_ce, "kl": loss_kl}
\end{lstlisting}

\textit{Соответствие \eqref{eq:totalloss}.} Параметры $\lambda_{\text{ce}}=0{,}05$, $\lambda_{\text{KL}}=0{,}1$, $\tau=2{,}0$ выбраны абляцией; масштаб $s=0{,}10$ выбран по результатам поиска веса в [07\_wav2lip\_finetune\_H1.ipynb](07\_wav2lip\_finetune\_H1.ipynb). Аудиологиты дистиллируются с \texttt{detach()}, поскольку аудиоэнкодер на этапе 5 заморожен и не должен принимать градиент. Косинусный член из \eqref{eq:totalloss} в финальной конфигурации обнуляется ($\lambda_{\cos}=0$) на основании отрицательного $\Delta$F1 наивной косинусной формулировки в абляции (см. табл.~\ref{tab:ablation}).

\newpage
\section{\centering}\label{PrilV}
\begin{center}
\hspace*{-7mm}{\large\textbf{Полные таблицы перебора энкодеров (этап 2)}}
\end{center}
\setcounter{equation}{0}
\renewcommand{\theequation}{\Asbuk{section}.\arabic{equation}}
\setcounter{table}{0}
\renewcommand{\thetable}{\Asbuk{section}.\arabic{table}}
\setcounter{figure}{0}
\renewcommand{\thefigure}{\Asbuk{section}.\arabic{figure}}

В рамках этапа 2 проведено 22 запуска (11 аудио, 11 видео); ниже приведены полные таблицы валидационной макро-F1 для всех конфигураций. Источник данных --- \texttt{02\_train\_emotion\_encoders.ipynb} (cell~12, RESULTS SUMMARY) и проект W\&B \texttt{uncanny-valley-encoders-4emo}.

\begin{table}[H]
\caption{Полный перебор аудиоэнкодеров эмоций (валидационная макро-F1)}
\label{tab:full_audio_sweep}
\centering
\small
\begin{tabular}{|p{6.5cm}|p{1.8cm}|p{1.8cm}|p{1.6cm}|}
\hline
\textbf{Конфигурация} & \textbf{Базовая модель} & \textbf{Темп обучения} & \textbf{Best Val F1} \\
\hline
4emo-w2v2-er-lr3e5 & wav2vec2-base-superb-er & $3\!\cdot\!10^{-5}$ & \textbf{0,8264} \\
\hline
4emo-hubert-lg-lr2e5 & hubert-large-ll60k & $2\!\cdot\!10^{-5}$ & 0,8036 \\
\hline
4emo-w2v2-lg-lr2e5 & wav2vec2-large & $2\!\cdot\!10^{-5}$ & 0,8011 \\
\hline
4emo-w2v2-er-lr5e5 & wav2vec2-base-superb-er & $5\!\cdot\!10^{-5}$ & 0,8010 \\
\hline
4emo-w2v2-er-lr1e4 & wav2vec2-base-superb-er & $1\!\cdot\!10^{-4}$ & 0,7965 \\
\hline
4emo-hubert-er-lr1e4 & hubert-base-superb-er & $1\!\cdot\!10^{-4}$ & 0,7682 \\
\hline
4emo-w2v2-er-lr5e5-w5 & wav2vec2-base-superb-er, window=5s & $5\!\cdot\!10^{-5}$ & 0,7582 \\
\hline
4emo-hubert-er-lr3e5-nf & hubert-base-superb-er, no-freeze & $3\!\cdot\!10^{-5}$ & 0,7536 \\
\hline
4emo-hubert-er-lr5e5 & hubert-base-superb-er & $5\!\cdot\!10^{-5}$ & 0,7449 \\
\hline
4emo-hubert-er-lr3e5-ls0 & hubert-base-superb-er, label\_smoothing=0 & $3\!\cdot\!10^{-5}$ & 0,7389 \\
\hline
4emo-hubert-er-lr3e5 & hubert-base-superb-er & $3\!\cdot\!10^{-5}$ & 0,7343 \\
\hline
\end{tabular}
\end{table}

\begin{table}[H]
\caption{Полный перебор видеоэнкодеров эмоций (валидационная макро-F1)}
\label{tab:full_video_sweep}
\centering
\small
\begin{tabular}{|p{6.5cm}|p{1.8cm}|p{1.8cm}|p{1.6cm}|}
\hline
\textbf{Конфигурация} & \textbf{Базовая модель} & \textbf{Темп обучения} & \textbf{Best Val F1} \\
\hline
4emo-tsf-lr3e5-8f & timesformer-base-finetuned-k400, 8 кадров & $3\!\cdot\!10^{-5}$ & \textbf{0,7898} \\
\hline
4emo-tsf-lr3e5-16f-nf & timesformer-base, 16 кадров, no-freeze & $3\!\cdot\!10^{-5}$ & 0,7805 \\
\hline
4emo-tsf-lr5e5-16f & timesformer-base, 16 кадров & $5\!\cdot\!10^{-5}$ & 0,7796 \\
\hline
4emo-tsf-lr3e5-16f-wd05 & timesformer-base, 16 кадров, wd=0,05 & $3\!\cdot\!10^{-5}$ & 0,7715 \\
\hline
4emo-tsf-lr1e5-16f & timesformer-base, 16 кадров & $1\!\cdot\!10^{-5}$ & 0,7711 \\
\hline
4emo-tsf-lr3e5-16f & timesformer-base, 16 кадров & $3\!\cdot\!10^{-5}$ & 0,7706 \\
\hline
4emo-tsf-lr3e5-16f-cb & timesformer-base, 16 кадров, class-balanced & $3\!\cdot\!10^{-5}$ & 0,7706 \\
\hline
4emo-tsf-lr3e5-16f-ls0 & timesformer-base, 16 кадров, label\_smoothing=0 & $3\!\cdot\!10^{-5}$ & 0,7686 \\
\hline
4emo-tsf-lr1e5-16f-f3 & timesformer-base, 16 кадров, freeze 3 layers & $1\!\cdot\!10^{-5}$ & 0,7629 \\
\hline
4emo-tsf-lr3e5-12f & timesformer-base, 12 кадров & $3\!\cdot\!10^{-5}$ & 0,7453 \\
\hline
4emo-tsf-lr3e5-4f & timesformer-base, 4 кадра & $3\!\cdot\!10^{-5}$ & 0,7092 \\
\hline
\end{tabular}
\end{table}

\textit{Анализ перебора.} Аудиоэнкодеры на основе wav2vec~2.0 систематически превосходят HuBERT-варианты на этой задаче, причём базовая модель \texttt{wav2vec2-base-superb-er}, дообученная на IEMOCAP-производном корпусе SUPERB-ER, превосходит более крупные общего назначения варианты \texttt{wav2vec2-large} и \texttt{hubert-large-ll60k}. Это согласуется с общим принципом доменной адаптации: pre-training на близком домене даёт б\'ольший эффект, чем масштабирование архитектуры. Для видео --- TimeSformer на 8 кадрах даёт наилучшую валидационную F1 при наименьших вычислительных затратах среди всех конфигураций по числу кадров, что и обусловило выбор 8-кадровой версии для финальных экспериментов.

\newpage
\section{\centering}\label{PrilG}
\begin{center}
\hspace*{-7mm}{\large\textbf{Полные таблицы бутстрап-CI и per-actor sign-test}}
\end{center}
\setcounter{equation}{0}
\renewcommand{\theequation}{\Asbuk{section}.\arabic{equation}}
\setcounter{table}{0}
\renewcommand{\thetable}{\Asbuk{section}.\arabic{table}}
\setcounter{figure}{0}
\renewcommand{\thefigure}{\Asbuk{section}.\arabic{figure}}

Полная сводка результатов бутстрап-анализа (10\,000 реплик с возвращением, перцентильный 95~\%-CI, парные выборки) приведена ниже. Источник: [11\_bootstrap\_analysis.ipynb](11\_bootstrap\_analysis.ipynb), cell 9. Дополнительно включён per-actor sign-test ($n_{\text{актёров}}=3$ на тестовом сплите).

\begin{table}[H]
\caption{Бутстрап-95~\% CI на $\Delta\mathrm{F1}$, val + test, четыре эвалюатора}
\label{tab:full_bootstrap}
\centering
\small
\begin{tabular}{|p{2cm}|p{1.7cm}|p{2.5cm}|c|c|c|}
\hline
\textbf{Эвалюатор} & \textbf{Split} & \textbf{Конфигурация} & $\Delta\mathrm{F1}$ & \textbf{95~\% CI} & \textbf{CI $\not\ni 0$} \\
\hline
Rajaram1996 & test & cekl-05 & $+0{,}089$ & $[+0{,}004,\,+0{,}176]$ & \textbf{да} \\
\hline
Rajaram1996 & val & cekl-04 & $+0{,}094$ & $[+0{,}024,\,+0{,}166]$ & \textbf{да} \\
\hline
Rajaram1996 & val & abl-ce-kl & $+0{,}072$ & $[+0{,}002,\,+0{,}141]$ & \textbf{да} \\
\hline
Rajaram1996 & val & cekl-01 & $+0{,}060$ & $[+0{,}002,\,+0{,}121]$ & \textbf{да} \\
\hline
Rajaram1996 & val & abl-cos-only & $+0{,}047$ & $[+0{,}012,\,+0{,}091]$ & \textbf{да} \\
\hline
motheecreator & val & abl-cos-only & $-0{,}046$ & $[-0{,}098,\,-0{,}004]$ & да (отриц.) \\
\hline
trpakov & test & abl-baseline & $-0{,}031$ & $[-0{,}068,\,-0{,}000]$ & да (отриц.) \\
\hline
\end{tabular}
\end{table}

Для всех остальных конфигураций бутстрап-CI на тесте включает ноль; полная выгрузка доступна в JSON-формате в репозитории под путём \texttt{11\_bootstrap\_analysis/bootstrap\_summary.csv}.

\begin{table}[H]
\caption{Per-actor sign-test для лучших конфигураций, тестовая выборка}
\label{tab:full_per_actor}
\centering
\small
\begin{tabular}{|p{2cm}|p{2.5cm}|c|c|c|c|}
\hline
\textbf{Эвалюатор} & \textbf{Конфигурация} & $n_+$ & $n_-$ & \textbf{Mean $\Delta$} & \textbf{Median $\Delta$} \\
\hline
Rajaram1996 & cekl-05 & 2 & 1 & $+0{,}063$ & $+0{,}094$ \\
\hline
Rajaram1996 & cekl-01 & 2 & 1 & $+0{,}010$ & $+0{,}031$ \\
\hline
Rajaram1996 & cekl-04 & 1 & 1 & $+0{,}010$ & $0{,}000$ \\
\hline
Rajaram1996 & abl-ce-cos & 2 & 1 & $+0{,}010$ & $+0{,}031$ \\
\hline
motheecreator & abl-ce-cos & 2 & 0 & $+0{,}104$ & $+0{,}125$ \\
\hline
motheecreator & abl-ce-kl & 2 & 0 & $+0{,}094$ & $+0{,}063$ \\
\hline
motheecreator & cekl-04 & 2 & 1 & $+0{,}094$ & $+0{,}094$ \\
\hline
motheecreator & abl-ce-only & 2 & 0 & $+0{,}094$ & $+0{,}063$ \\
\hline
motheecreator & cekl-05 & 2 & 1 & $+0{,}073$ & $+0{,}063$ \\
\hline
motheecreator & cekl-01 & 3 & 0 & $+0{,}063$ & $+0{,}063$ \\
\hline
dima806 & abl-ce-cos & 1 & 1 & $+0{,}052$ & $0{,}000$ \\
\hline
dima806 & abl-ce-only & 1 & 1 & $+0{,}052$ & $0{,}000$ \\
\hline
dima806 & cekl-04 & 2 & 1 & $+0{,}042$ & $+0{,}031$ \\
\hline
\end{tabular}
\end{table}

\newpage
\section{\centering}\label{PrilD}
\begin{center}
\hspace*{-7mm}{\large\textbf{Поэмоциональные матрицы ошибок и оценки}}
\end{center}
\setcounter{equation}{0}
\renewcommand{\theequation}{\Asbuk{section}.\arabic{equation}}
\setcounter{table}{0}
\renewcommand{\thetable}{\Asbuk{section}.\arabic{table}}
\setcounter{figure}{0}
\renewcommand{\thefigure}{\Asbuk{section}.\arabic{figure}}

Ниже приведены матрицы ошибок и поэмоциональные оценки P/R/F1 для отобранных моделей и эвалюаторов на тестовой выборке. Источник: \texttt{08\_wav2lip\_external\_evaluation.ipynb}, cell 16, артефакты \texttt{cm\_*.png}.

\begin{table}[H]
\caption{Матрица ошибок baseline на motheecreator/test (rows: true, cols: pred)}
\label{tab:cm_baseline_mothee}
\centering
\small
\begin{tabular}{|c|c|c|c|c|}
\hline
true / pred & happy & sad & angry & disgust \\
\hline
happy & 24 & 0 & 0 & 0 \\
\hline
sad & 13 & 3 & 0 & 8 \\
\hline
angry & 8 & 7 & 1 & 8 \\
\hline
disgust & 9 & 6 & 0 & 9 \\
\hline
\end{tabular}
\end{table}

\begin{table}[H]
\caption{Матрица ошибок abl-ce-cos на motheecreator/test (rows: true, cols: pred)}
\label{tab:cm_cecos_mothee}
\centering
\small
\begin{tabular}{|c|c|c|c|c|}
\hline
true / pred & happy & sad & angry & disgust \\
\hline
happy & 21 & 3 & 0 & 0 \\
\hline
sad & 1 & 23 & 0 & 0 \\
\hline
angry & 6 & 17 & 0 & 1 \\
\hline
disgust & 1 & 20 & 0 & 3 \\
\hline
\end{tabular}
\end{table}

\textit{Качественное наблюдение.} На motheecreator после CE+COS-файнтюнинга baseline-смещение в сторону <<happy>> (по 24 ошибки на классы sad/angry/disgust в столбце <<happy>> у baseline) частично заменяется на смещение в сторону <<sad>> (17--20 ошибок в столбце <<sad>> для angry/disgust у abl-ce-cos). Поэтому общая макро-F1 растёт, но per-class F1 для \textit{angry} становится 0 из-за полного провала precision для этого класса на данном эвалюаторе. Это иллюстрирует ограничение макро-F1 как агрегата: положительная динамика по большинству классов может сопровождаться катастрофой на одном классе. Межклассификаторное согласие смягчает данный риск через требование положительного $\Delta$F1 одновременно на двух эвалюаторах.

\begin{table}[H]
\caption{Матрица ошибок baseline на Rajaram1996/test (rows: true, cols: pred)}
\label{tab:cm_baseline_raja}
\centering
\small
\begin{tabular}{|c|c|c|c|c|}
\hline
true / pred & happy & sad & angry & disgust \\
\hline
happy & 9 & 0 & 6 & 9 \\
\hline
sad & 0 & 0 & 19 & 5 \\
\hline
angry & 2 & 1 & 17 & 4 \\
\hline
disgust & 1 & 0 & 12 & 11 \\
\hline
\end{tabular}
\end{table}

\begin{table}[H]
\caption{Матрица ошибок cekl-05 на Rajaram1996/test (rows: true, cols: pred)}
\label{tab:cm_cekl05_raja}
\centering
\small
\begin{tabular}{|c|c|c|c|c|}
\hline
true / pred & happy & sad & angry & disgust \\
\hline
happy & 8 & 2 & 4 & 10 \\
\hline
sad & 0 & 18 & 3 & 3 \\
\hline
angry & 3 & 13 & 4 & 4 \\
\hline
disgust & 1 & 5 & 5 & 13 \\
\hline
\end{tabular}
\end{table}

\textit{Качественное наблюдение.} На Rajaram1996 baseline имеет default-смещение к <<angry>>; cekl-05-файнтюнинг ослабляет это смещение и даёт значительно лучший recall на \textit{sad} (с 0/24 = 0,00 до 18/24 = 0,75) и \textit{disgust} (с 11/24 = 0,46 до 13/24 = 0,54). Это объясняет, почему именно на Rajaram1996 cekl-05 проходит бутстрап-критерий значимости.

\newpage
\section{\centering}\label{PrilE}
\begin{center}
\hspace*{-7mm}{\large\textbf{Сводная статистическая таблица всех тестов}}
\end{center}
\setcounter{equation}{0}
\renewcommand{\theequation}{\Asbuk{section}.\arabic{equation}}
\setcounter{table}{0}
\renewcommand{\thetable}{\Asbuk{section}.\arabic{table}}
\setcounter{figure}{0}
\renewcommand{\thefigure}{\Asbuk{section}.\arabic{figure}}

\begin{table}[H]
\caption{Сводка статистических тестов для финальных моделей Wav2Lip (test, $n=96$)}
\label{tab:full_stats}
\centering
\small
\begin{tabular}{|p{3.5cm}|p{4cm}|c|c|}
\hline
\textbf{Метрика} & \textbf{Тест} & \textbf{Статистика} & $p$-value \\
\hline
Internal F1 (cekl-01 vs baseline) & McNemar ($n_{10}=17, n_{01}=9$) & $\chi^2 = 1{,}88$ & 0,170 \\
\hline
Internal F1 (cekl-01 vs baseline) & Bootstrap 95~\% CI & $[-0{,}03,\,+0{,}21]$ & --- \\
\hline
External F1 motheecreator (cekl-01 vs baseline) & McNemar ($n_{10}=8, n_{01}=14$) & $\chi^2 = 1{,}14$ & 0,286 \\
\hline
External F1 motheecreator (cekl-01 vs baseline) & Bootstrap 95~\% CI & $[-0{,}05,\,+0{,}14]$ & --- \\
\hline
External F1 Rajaram1996 (cekl-05 vs baseline) & McNemar ($n_{10}=14, n_{01}=21$) & $\chi^2 = 1{,}03$ & 0,310 \\
\hline
External F1 Rajaram1996 (cekl-05 vs baseline) & Bootstrap 95~\% CI & $[+0{,}004,\,+0{,}176]$ & \textbf{знач.} \\
\hline
LSE-C (cekl-01 vs baseline) & Парный $t$-критерий & $t = 0{,}51$ & 0,611 \\
\hline
L1 reconstruction (cekl-01 vs baseline) & Парный $t$-критерий & $t = -5{,}31$ & $7{,}3\!\cdot\!10^{-7}$ \\
\hline
L1 reconstruction (cekl-01 vs baseline) & Уилкоксон & $W = 346$ & $4{,}4\!\cdot\!10^{-13}$ \\
\hline
Spearman $\rho$ (int\_val\_F1, ext\_test\_F1\_mothee) & Корреляция Спирмена & $\rho = +0{,}612$ & 0,06 \\
\hline
\end{tabular}
\end{table}

\begin{table}[H]
\caption{Сводка статистических тестов для SadTalker (val coefficient + test rendered)}
\label{tab:full_stats_sad}
\centering
\small
\begin{tabular}{|p{3.5cm}|p{4cm}|c|c|}
\hline
\textbf{Метрика} & \textbf{Тест} & \textbf{Статистика} & $p$-value \\
\hline
Internal coef-level F1 (abl-ce-cos vs baseline, val) & McNemar & $\chi^2 = 6{,}72$ & \textbf{0,0095} \\
\hline
Internal coef-level F1 (abl-ce-only vs baseline, val) & McNemar & $\chi^2 = 3{,}84$ & 0,050 \\
\hline
External motheecreator F1 rendered (abl-ce-cos vs baseline, $n=24$) & McNemar & $\chi^2 = 0{,}25$ & 0,617 \\
\hline
External Rajaram1996 F1 rendered (abl-ce-cos vs baseline, $n=24$) & McNemar & $\chi^2 = 2{,}25$ & 0,134 \\
\hline
External trpakov F1 rendered (cekl-025 vs baseline) & McNemar & $\chi^2 = 0{,}00$ & 1,000 \\
\hline
\end{tabular}
\end{table}

Совокупность представленных в Приложениях А--Д таблиц покрывает все эмпирические утверждения, сделанные в основном тексте, и обеспечивает прозрачность для воспроизведения результатов внешним исследователем.
%\captionof{lstlisting}{Фрагмент кода функции subfinder}  % ✅ Русская подпись отдельно
% \begin{minted}[frame=lines, framesep=2mm, fontsize=\small, label=some-code,caption=Фрагмент кода функции subfinder]{python}
% def subfinder(words_list, answer_list): # code of function
%   matches = []
%   start_indices = []
%   end_indices = []
%   for idx, i in enumerate(range(len(words_list))):
%     if words_list[i] == answer_list[0] and words_list[i : i + len(answer_list)] == answer_list:
%       matches.append(answer_list)
%       start_indices.append(idx)
%       end_indices.append(idx + len(answer_list) - 1)
%     if matches:
%       return matches[0], start_indices[0], end_indices[0]
%     else:
%       return None, 0, 0
% \end{minted}

% Импорт из внешнего файла:
% \lstinputlisting[language=Python, frame=single]{script.py}



% \intro{Название приложения}\label{PrilA}


% Содержание приложения

% \newpage
% \intro{Название приложения}
% \label{PrilB}

% Содержание приложения

\end{document}

